% Modernised reading — fonds Alexandre Grothendieck, shelfmark 105, pages 1-52
% (the whole folder).
%
% This is an INTERPRETATION, not a transcription. It re-reads the pages in
% today's notation and vocabulary, states the hypotheses the manuscript leaves
% implicit, and drops the critical apparatus so that the mathematics reads
% straight through. Everything the brackets and underlines carried moves into
% a footnote. It is held to being mathematically correct as it stands; where
% the manuscript is loose or mistaken, the true statement is what appears here
% and the footnote says what the page has. In any dispute of substance the
% transcription governs: whoever cites Grothendieck must cite
% batch-0{1,2,3}.fr.tex.
%
% Scope: the folder was read whole — three batches, pages 1-52, thirty pages
% carrying his hand — before any of this was written, and it is ONE file for
% the shelfmark.
%
% What the folder is. Two French scratch leaves (pp. 1-2); an English run
% paginated by him 1-16 (pp. 4-39) on model structures in a topos with the
% monomorphisms as cofibrations, in three stretches: generation from a small
% set (TC)_0 and the comparison of W~ = TF o TC with W (pp. 4-19, his 1-8);
% a restart from the pair (C, W), with saturated classes, weak factorization
% systems and a criterion for a model structure (pp. 21-31, his 9-12); the
% small-object argument that supplies the factorization both stretches had
% admitted (pp. 33-39, his 13-16). Then a French run paginated 6-18
% (pp. 42-52), whose pages 1-5 are not in the folder, on axioms W7, W8 for a
% class W in Cat tested on W_D, the cohomological equivalences of a
% dérivateur D, ending on the axiom « Der 6 »: a square is cartesian iff
% cocartesian.
%
% Notation fixed for the whole document.
%  - The topos is written \mathcal{E} (he writes A, which is also his letter
%    for the source of an arrow i : A -> B); e its final object.
%  - Phi^* (left lifting) and Phi_* (right lifting) are his, from p. 24, used
%    from p. 4 on; i \pitchfork f means i has the LLP with respect to f.
%  - His underlined Hom (internal) is \mathbf{Hom}; his boxed cross is the
%    pushout-product \boxtimes; alpha(i,f) the pullback-hom.
%  - His « Quillen pair » is written « système de factorisation faible », his
%    « Quillen triple » « structure de modèles »: the modern « paire de
%    Quillen » is an adjunction, and the clash is footnoted.
%  - The section of an interval is \varepsilon : e -> I (he writes e for both).
%  - His underlined W (a localisateur fondamental on Cat) is \mathcal{W}; the
%    test category of p. 19 b) is B (he writes T, also the object of p. 6).
%  - Part 3: Psi = {0 < a, 0 < b}, Q = Delta^1 x Delta^1, k : Psi -> Q,
%    Phi = Psi^op = (a -> 0 <- b) (he calls its apex 1 on p. 49); \mathcal{A}
%    = D(e); the Grothendieck construction s |-> X_{/s} is \mathcal{G}_I(X)
%    (he writes a thick Psi with an index, which collides with the poset Psi);
%    objects of D(I) are read as diagrams of type I^op, as p. 45 says.
%
% Substantive departures from the pages, each footnoted where it occurs:
% p. 2, the composites s o p, b o p are a fibred and a cofibred functor
% (the page seems to deny it; reading uncertain); p. 6, the NB's claim of
% what is used is sharpened (only retracts); p. 9, (iii) is equivalent to
% boxtimes-stability, not to stability under ? x T; p. 11, fourth line of
% the table is C boxtimes TC (page hatched), and « F stable » is TF; pp. 11-13,
% the passage from generators to TC boxtimes C is supplied; p. 13, his M1/M5
% are not Quillen's numbering; p. 15, TF in W from axiom (1) needs c) of p. 4,
% argument supplied; p. 17, the factorization used must be the cellular one;
% p. 21, his 3), 4), 5) are (3)(4)(5) of pp. 15-17; p. 26, the question is
% answered (no) by a counterexample of ours; p. 27, factorization implicit
% in (i); p. 37, (d) for filtering limits is more than the Zorn argument
% gives, p. 39 restricts it to well-ordered systems; p. 35, the « pi = 0 » is
% aleph_0 and alpha_0 a regular cardinal; p. 45, « i_0^* iso <=> i^* iso »
% holds only => for a cartesian square, as p. 46 itself says; p. 49,
% psi_!(xi) lives in D(Q^op); p. 50, the implications listed are the
% cocartesian ones; p. 52, the shape Phi normalised to (a -> 0 <- b).
% Remarks of ours, footnoted as such: W7 bis and W8 bis fail for W_infinity
% (acyclic spaces) and hold for homology equivalences, which is what makes
% « Der 6 » the right axiom; Der 8's a + b => c holds in the stable case;
% the statements « admitted in HOT » hold for connective complexes and fail
% for homotopy types of spaces.
%
% Announced and not established: the factorization (admitted pp. 13, 21,
% supplied pp. 33-39 only under the smallness of Phi_0, which p. 39 wants to
% remove); W ⊂ W~ (p. 19, « [ ? ??? ] »), and the two cases a), b) wanted
% there; why TC = C ∩ W is the saturation of (TC)_0 (p. 31); the « à
% vérifier » of p. 45; Der 6 is posed, not derived; the lemma of p. 50 is
% left at « forme de Der 3 »; Der 8's « (?) »; W7 bis, W8 bis are stated,
% not proved; the struck « cqfd » of p. 42 closes an argument whose start
% (his pp. 1-5) is not in the folder.
%
% Pass: Opus 5.5 (claude-opus-5-5), 2026-10-03 — first pass, unchecked against the pages by a human.
\documentclass[11pt,a4paper]{article}
\input{../preamble/grothendieck}

\folder{105}
\pages{1}{52}
\dating{[à partir de 1982]}
\watermark{Édition de démonstration\\interprétation personnelle de l'œuvre}
\foldertitle{[Champs (stacks) 2] : notes manuscrites (s.d.) — lecture modernisée du dossier entier}

\begin{document}

\section*{Résumé}

\begin{resume}
En topologie, deux formes ont « le même type d'homotopie » quand on peut
déformer l'une en l'autre sans rien déchirer : un disque et un point, un
anneau et un cercle. Pour calculer avec cette notion, il faut savoir quelles
applications on déclare être des \emph{équivalences}, et cela ne suffit pas :
il faut encore deux familles d'applications commodes, de bonnes inclusions et
de bonnes projections, qui se relèvent les unes par rapport aux autres comme un
chemin se relève dans un revêtement. Quillen a appelé en 1967 la donnée de ces
trois familles une structure de modèles fermée. Ces feuillets, en anglais pour
la plupart, cherchent à quelles conditions une classe d'équivalences donnée
d'avance, dans un topos, est celle d'une telle structure dont les bonnes
inclusions sont toutes les inclusions.

La question vient de \emph{Pursuing Stacks}, auquel l'inventaire rattache le
dossier. Les types d'homotopie y sont décrits par des préfaisceaux sur des
« catégories test », et les équivalences y sont définies tout de suite, par un
critère homologique. Chez Kan et Quillen, au contraire, on part d'un petit
ensemble d'inclusions élémentaires, les cornets des ensembles simpliciaux, et
les équivalences sont ce qui en sort. Grothendieck le dit lui-même à la page
31 : de son point de vue, l'équivalence est donnée d'avance, « in a rather
tangible way », et c'est la construction de Kan et Quillen qui lui paraît ad hoc.
Le dossier cherche le pont entre les deux. Quand les inclusions qui sont des
équivalences sont-elles exactement celles qu'on obtient en saturant quelques
inclusions élémentaires ?

L'idée est celle d'un intervalle. On se donne quelques objets $I$ munis d'un
point, qu'on veut voir comme des segments ; on fait glisser une inclusion
$A_0 \subset B_0$ le long de $I$, ce qui donne une inclusion de « cylindres » ;
et l'on prend pour inclusions élémentaires toutes celles-là. Le segment
universel d'un topos est l'objet de Lawvere $\Omega$, qui classe les
sous-objets. Il a toujours les propriétés voulues, ce qui arrache à la marge de
la page 15 trois points d'exclamation. Reste à fabriquer les factorisations,
et c'est l'argument par lequel on plonge un module dans un module injectif : on
ajoute toutes les solutions qui manquent, puis on recommence, transfiniment,
jusqu'à ce que plus rien ne manque. Les pages 33 à 39 l'écrivent en entier ;
c'est l'« argument du petit objet ».

Le dossier procède comme on le fait en travaillant. Il admet d'abord les
factorisations, pour le moment, pour voir ce qu'il reste à démontrer. Il
reprend ensuite tout depuis l'autre bout, sous le titre « Alternative
presentation of theory ». Il ne revient aux factorisations qu'une fois le reste
en place. Les deux présentations ne sont pas fondues : la première engendre,
la seconde part de ce qu'on veut obtenir.

Les onze dernières pages, en français, changent de cadre. Une catégorie y est
recouverte par deux morceaux ouverts, comme un espace par deux ouverts dans la
suite exacte de Mayer--Vietoris, et l'on teste deux axiomes de recollement sur
les équivalences que voit une cohomologie. Elles aboutissent à un axiome
qu'elles jugent provisoire : un carré y est « cartésien » (un produit fibré)
si et seulement s'il est « cocartésien » (une somme amalgamée). C'est le monde
des suites exactes longues, où recoller et intersecter reviennent au même ; ce
n'est pas celui des espaces.

Les noms modernes à chercher : systèmes de factorisation faibles, catégories de
modèles à engendrement cofibrant, structures de modèles de Cisinski sur les
topos, extensions anodines ; et, pour la fin, dérivateurs stables et limites
inductives homotopiques de Thomason.

\keywords{lifting property, weak factorization system, model category, pushout-product, monoidal model category, subobject classifier, interval object, anodyne extension, Cisinski model structure, test category, basic localizer, small object argument, cofibrantly generated model category, accessible category, proper model category, 2-out-of-6 property, Thomason model structure, derivator, homotopy cartesian square, stable derivator, Grothendieck construction, homotopy pushout, Mayer–Vietoris sequence, homology equivalence}
\end{resume}

\section*{Le fil du dossier, et les conventions}

Le dossier a cinquante-deux pages ; trente portent sa main. Les autres sont des
feuilles de listing d'ordinateur (journaux de travaux, pages de garde,
tableaux) qui ne portent rien de lui, des feuilles blanches, un carton rose
(page 3) où ne figure qu'un nom biffé, et une page dactylographiée (page 43,
numérotée « 90 » par la dactylo) d'une prose sans rapport avec les
mathématiques. Plusieurs des feuillets écrits le sont au dos de listings ; celui
de la page 27 porte imprimée la date « 06/09/82 », qui borne ce feuillet par
en dessous et lui seul.

\begin{enumerate}
\item[1.] \emph{Deux brouillons} (pages 1 et 2, en français, non paginés) : un
lemme de relèvement dans $\mathrm{Cat}$, et la question d'une structure de
Quillen sur $\mathrm{Cat}$.
\item[2.] \emph{Une suite en anglais, paginée 1 à 16 de sa main} (pages 4 à
39 ; sa page « 10 bis » est la page 27, et la page 26 est un brouillon sans
numéro). Elle contient trois étapes, 2a à 2c.
\item[2a.] Pages 4 à 19 (ses pages 1 à 8) : d'un petit ensemble $(TC)_0$
d'inclusions élémentaires on tire quatre classes $C, TC, F, TF$ et une classe
$\widetilde W = TF \circ TC$. Le produit-boîte, puis l'exemple des cylindres
sur des intervalles, montrent la compatibilité au produit. Le programme est
$\widetilde W = W$ : l'inclusion $\widetilde W \subset W$ s'obtient sous cinq
axiomes, l'inclusion inverse reste ouverte.
\item[2b.] Pages 21 à 31 (ses pages 9 à 12) : la reprise. On part de la paire
$(C, W)$ et l'on pose directement $TC = C \cap W$. Viennent les classes
saturées, les systèmes de factorisation faibles, la définition d'une structure
de modèles par deux tels systèmes, et un critère d'existence (page 29).
\item[2c.] Pages 33 à 39 (ses pages 13 à 16), « Factorization » : l'argument
du petit objet, qui fournit la factorisation admise en 2a et 2b.
\item[3.] \emph{Une suite en français, paginée 6 à 18 de sa main} (pages 42 à
52), dont les pages 1 à 5 ne sont pas dans le dossier : les axiomes W7, W8 et
leurs formes « bis » pour une classe $W$ de flèches de $\mathrm{Cat}$, testés
sur la classe $W_{\mathbb D}$ des équivalences cohomologiques d'un dérivateur
$\mathbb D$, et l'axiome « Der 6 ».
\end{enumerate}

Rien ne date la troisième suite par rapport à la deuxième. Son vocabulaire
(dérivateur, axiomes « Der ») est celui du long manuscrit de Grothendieck sur
les dérivateurs ; c'est un rapprochement, non une datation.

\emph{Conventions.} Pour une catégorie $M$, $\mathrm{Fl}(M)$ est l'ensemble de
ses flèches. On écrit $i \pitchfork f$ quand $i$ a la propriété de relèvement à
gauche par rapport à $f$ (tout carré de $i$ vers $f$ admet une diagonale), et,
pour $\Phi \subset \mathrm{Fl}(M)$,
\[
\Phi^* = \{u \mid u \pitchfork \varphi \ \ \forall \varphi \in \Phi\}, \qquad
\Phi_* = \{u \mid \varphi \pitchfork u \ \ \forall \varphi \in \Phi\},
\]
notation qui est la sienne à partir de la page 24 et qu'on emploie dès la page
4. Le topos des pages 4 à 19 est noté $\mathcal{E}$, son objet final
$e$.\footnote{Il écrit $A$ pour le topos, et aussi pour la source d'une
flèche $i : A \to B$ ; on lève la collision.} Le $\mathbf{Hom}$ en gras est le
$\mathrm{Hom}$ interne, qu'il souligne. Les chiffres et les lettres qu'il entoure (ses axiomes
numérotés de la page 15 à la page 17, ses conditions des pages 21 à 39) sont
rendus entre parenthèses, (1) à (5) et (a) à (e).

Trois mots changent de sens d'une partie à l'autre, et on les fixe ici. Une
\emph{fibration}, sans autre précision, est une flèche de la classe $F$ d'une
structure de modèles ; une \emph{catégorie fibrée} ou \emph{cofibrée} sur une
base est prise au sens de Grothendieck (SGA~1), et ce sont elles seules qui
paraissent dans la troisième partie. Sa « paire de Quillen » est ici un
\emph{système de factorisation faible}, et son « triple de Quillen » une
\emph{structure de modèles}. Pour la saturation d'une classe $W$, enfin, on
distingue trois degrés. $W$ est \emph{faiblement saturée} si elle contient les
isomorphismes, a la propriété du deux-sur-trois et la propriété c) de la page
4 ; il dit aussi « mildly ». Elle est \emph{saturée} si $fg, gf \in W$
entraînent $f, g \in W$ (page 26). Elle est \emph{fortement saturée} si
$f \in W$ équivaut à ce que l'image de $f$ dans $W^{-1}M$ soit un
isomorphisme.

Les notations de la troisième partie sont fixées à son entrée (page 42).

\emph{Le dossier annonce et n'établit pas} : les factorisations, admises aux
pages 13 et 21, puis obtenues aux pages 33 à 39 sous une hypothèse de
petitesse dont la page 39 voudrait se défaire ; l'inclusion
$W \subset \widetilde W$ (page 19) et les deux cas où il la voudrait ; la
raison pour laquelle $C \cap W$ serait la saturation de $(TC)_0$ (page 31) ;
l'axiome « Der 6 », posé et non déduit ; le lemme de la page 50, arrêté sur
« forme de Der 3 » ; les énoncés W7 bis et W8 bis, posés sans démonstration.

\pagerange{1}{2}
\section*{Deux brouillons (pages 1 et 2)}

\emph{Page 1 : un lemme de relèvement dans $\mathrm{Cat}$.} Soit
$\partial \to I$ l'inclusion des extrémités d'un intervalle.\footnote{Sous
$I$, une mention entre parenthèses illisible ; plus bas, deux lignes biffées,
dont « $\Pi_0(\partial) \geq 2$ », qui semble demander que les extrémités
soient au moins deux. En haut à gauche, un signe en forme de Z barré, suivi de
« a », non identifié.} On considère les flèches
\[
\varnothing \to e, \qquad \partial \times S^n(N) \longrightarrow I \times S^n(N),
\]
et des « cônes » construits sur elles, parmi lesquels
$\{0\} \to \Delta^1$ et $\{1\} \to \Delta^1$.\footnote{Les objets $S^n(N)$ ne
sont définis nulle part sur les feuillets ; $N$ parcourt un ensemble
$\mathcal N$ de parties, lues « finies » avec doute, de $\mathbb N$. Le mot qui
suit « cônes » est illisible, et l'abréviation « compl. droits » n'est pas
résolue. Un $\mathbf{Hom}(T, X)$ isolé est écrit au milieu de la page ; la
lettre, très appuyée, pourrait aussi être le $I$ du haut.} Le lemme énonce :
\begin{enumerate}
\item[a)] si $X \to e$ a la propriété de relèvement à droite par rapport à
$\varnothing \to e$ et aux $\partial \times S^n(N) \to I \times S^n(N)$, alors
$X$ est $W_\infty$-asphérique, c'est-à-dire que son nerf est contractile ;
\item[b)] si $X \to Y$ a la propriété de relèvement à droite par rapport aux
cônes sur ces flèches, alors $X \to Y$ est $W_\infty$-propre (resp.~$W_\infty$-lisse).
\end{enumerate}
C'est l'analogue, dans $\mathrm{Cat}$, du fait qu'un ensemble simplicial ayant
la propriété de relèvement par rapport aux inclusions $\partial\Delta^n \subset
\Delta^n$ est contractile.\footnote{Le rapprochement est le nôtre. $W_\infty$
est la classe des foncteurs dont le nerf est une équivalence faible
d'homotopie, le plus petit localisateur fondamental (D.-C.~Cisinski,
2004). Les notions de foncteur $W$-propre et $W$-lisse sont celles de
\emph{Pursuing Stacks}, calquées sur les théorèmes de changement de base
propre et lisse de SGA~4 ; « lisse » est lu avec doute, aux trois endroits où
le mot paraît.} Le corollaire qui suit, très surchargé et en partie couvert
d'une tache d'encre, ne se lit que par fragments : une propriété de relèvement
à droite rendrait $X \to Y$ propre à fibres asphériques, et, sous une condition
illisible portant sur des catégories finies, « $W_\infty$-parfait » (lu avec
doute).\footnote{Le $Y$ de $X \to Y$ est écrit sur un $e$ ; plusieurs mots
biffés se superposent au texte. On ne restitue pas le corollaire.}

\emph{Page 2 : la flèche $\mathrm{Ar}(C) \to C \times C$.} Le feuillet est barré
de deux longs traits obliques et se lit sous la rature. Pour une catégorie $C$
et sa catégorie des flèches $\mathrm{Ar}(C)$, on voudrait que
\[
p = (s, b) : \mathrm{Ar}(C) \longrightarrow C \times C,
\]
qui envoie une flèche sur sa source et son but, ait des propriétés de type
fibration. Or $p$ n'est ni une catégorie fibrée ni une catégorie
cofibrée.\footnote{Elle est ce qu'on appelle aujourd'hui une fibration
discrète à deux faces (R.~Street, 1974) : fibrée en la première variable,
cofibrée en la seconde. Le nom n'est pas sur la page. La page ajoute que les
composées $s \circ p$, $b \circ p : \mathrm{Ar}(C) \to C$ « ne le sont pas » ;
le mot qui précède est illisible et « composées » est lu avec doute. Or
$s \circ p$ est fibrée (la précomposition $x \mapsto x u$ fournit les
flèches cartésiennes) et $b \circ p$ est cofibrée (par la postcomposition).
Chacune n'a que l'une des deux propriétés, et c'est le seul sens où la phrase
est juste ; on ne lit pas assez la page pour dire si c'est celui qu'elle a.}
Elle n'est pas non plus cohomologiquement propre ni lisse, et la page se
demande si au moins $s \circ p$ et $b \circ p$ le sont.\footnote{Avec la
convention que le dossier suit aux pages 44 et 52 (une catégorie cofibrée est
propre, une catégorie fibrée est lisse), $s \circ p$ est lisse et $b \circ p$
propre. La remarque est la nôtre.} « Mais ce sont plutôt les propriétés de
$p$ lui-même qui importent. » Il souhaiterait que $p$ soit une
« $W_\infty$-fibration » (« fibration cohomologique », lu avec doute), et
pose :

\textbf{Question.} Les $W_\infty$-fibrations et les cofibrations définies par
relèvement à gauche par rapport à elles forment-elles une structure de
modèles sur $\mathrm{Cat}$ ?\footnote{Une structure de modèles sur
$\mathrm{Cat}$ dont les équivalences faibles sont $W_\infty$ existe : c'est
celle de R.~Thomason (1980), obtenue par transfert le long de
$\mathrm{Ex}^2 N$. Ses fibrations sont définies par ce transfert et non de
façon intrinsèque, et la page ne définit pas les « $W_\infty$-fibrations ».
On ne dit donc pas si les deux classes coïncident, ni si $p$ est une
fibration de Thomason.}

\pagerange{4}{11}
\section*{Paires de relèvement dans un topos (pages 4 à 11)}

À partir d'ici le texte est en anglais, paginé de sa main à la tête de chaque
feuillet : la page 4 est sa page 1, la page 19 sa page 8, et la page 21 sa
page 9.

\subsection*{Les quatre classes}

Soit $\mathcal{E}$ un topos, et $W \subset \mathrm{Fl}(\mathcal{E})$ une
classe faiblement saturée :
\begin{enumerate}
\item[a)] les isomorphismes sont dans $W$ ;
\item[b)] si $f = vu$ et si deux des trois flèches $u, v, f$ sont dans $W$, la
troisième l'est ;
\item[c)] si $f : X \to Y$ et $g : Y \to X$ vérifient $gf = \mathrm{id}_X$ et
$fg \in W$, alors $f, g \in W$ ;
\end{enumerate}
« plus other conditions we are going to list, as we need them ».\footnote{« weakly » est ajouté au-dessus de la ligne. La fin de la
ligne b) mord sur le bord du feuillet. La condition c) est un cas particulier
de la stabilité par rétractes, qu'il posera page 17 (axiome (5)) ; la page 26 le
montre.}

Donnons-nous deux classes $(TC)_0 \subset C_0 \subset
\mathrm{Fl}(\mathcal{E})$, sans autre hypothèse pour l'instant, et
posons
\[
F = ((TC)_0)_*, \quad TF = (C_0)_*, \qquad TC = F^*, \quad C = (TF)^*.
\]
Comme $(TC)_0 \subset C_0$, on a $TF \subset F$, d'où $TC \subset C$, et
$(TC)_0 \subset TC$, $C_0 \subset C$. Chacune des deux paires $(TC, F)$ et
$(C, TF)$ se détermine par relèvement : $F = (TC)_*$ et $TF = C_*$, parce
que $\Phi_* = ((\Phi_*)^*)_*$ pour toute classe $\Phi$. Les classes $F$ et $TF$
sont stables par composition, changement de base et rétractes ; $C$ et $TC$
par composition, changement de cobase et rétractes. En marge : $C$ seront les
« cofibrations », $TC$ les « cofibrations triviales », $F$ les
« fibrations », $TF$ les « fibrations triviales ».

\textbf{Lemme} (page 6). Soit
$\widetilde W = \{ pi \mid p \in TF,\ i \in TC \}$. Alors
\[
TC = C \cap \widetilde W, \qquad TF = F \cap \widetilde W .
\]
En effet, $TC \subset C \cap \widetilde W$ en écrivant $i = \mathrm{id} \circ
i$. Inversement, soit $f = pi \in C$, $p \in TF$, $i \in TC$ ; comme
$f \pitchfork p$, le carré formé de $i$ et de l'identité de la cible a une
diagonale $s$, avec $sf = i$ et $ps = \mathrm{id}$, de sorte que $f$ est un
rétracte de $i$ et appartient à $TC$. L'énoncé sur $TF$ est dual.\footnote{La
page énonce les deux égalités sans démonstration ; celle-ci est la nôtre. Le
\emph{NB} qui suit dit que ces faits n'utilisent que l'existence de sommes
amalgamées et de produits fibrés dans $\mathcal{E}$, pas le fait que ce soit un
topos, ni $W$, ni la manière dont les paires ont été obtenues, seulement leur
relation symétrique et $TC \subset C$. C'est juste, et même plus que juste :
l'argument n'utilise que la stabilité par rétractes, c'est-à-dire aucune
limite.}

\subsection*{La flèche $\alpha(i,f)$ et le produit-boîte}

Pour $i : A \to B$ et $f : X \to Y$ dans $\mathcal{E}$, soit
\[
\alpha(i,f) : \mathbf{Hom}(B, X) \longrightarrow \mathbf{Hom}(i,f)
= \mathbf{Hom}(B, Y) \times_{\mathbf{Hom}(A, Y)} \mathbf{Hom}(A, X),
\]
qu'on appelle aujourd'hui la puissance fibrée de $f$ par $i$.\footnote{Sous
$i$ et $f$, biffés : « $i \in C$ », « $f \in F$ ». Le nom est le nôtre ;
on parle aussi de construction de Leibniz.} La page considère, pour un objet
$T$, les trois propriétés :
\begin{enumerate}
\item[(i)] $C$ et $TC$ sont stables par $? \times T$ ;
\item[(ii)] $TF$ et $F$ sont stables par $\mathbf{Hom}(T, ?)$ ;
\item[(iii)] si $i \in C$ et $f \in F$, alors $\alpha(i,f) \in F$, et
$\alpha(i,f) \in TF$ si de plus $i \in TC$ ou $f \in TF$.
\end{enumerate}

\textbf{(i) $\Leftrightarrow$ (ii)}, et plus précisément
$C \times T \subset C \Leftrightarrow \mathbf{Hom}(T, TF) \subset TF$ et
$TC \times T \subset TC \Leftrightarrow \mathbf{Hom}(T, F) \subset F$. Soit
en effet $f \in F$. Alors $\mathbf{Hom}(T, f) \in F = (TC)_*$ si et seulement
si $i \pitchfork \mathbf{Hom}(T, f)$ pour tout $i \in TC$, c'est-à-dire
$i \times \mathrm{id}_T \pitchfork f$, ce qui a lieu si $TC \times T \subset
TC$. La réciproque et l'autre paire se traitent de même (page 9).

Le \emph{NB} de la page 9 dégage ce qui sert. Soient $\otimes$ et
$\mathbf{Hom}$ deux foncteurs liés par l'isomorphisme d'adjonction
$\mathrm{Hom}(A \otimes T, X) \simeq \mathrm{Hom}(A, \mathbf{Hom}(T,
X))$.\footnote{Il l'appelle « relation de Cartan » ; c'est l'associativité
adjointe du livre de Cartan et Eilenberg.} Pour $j : A' \to B'$ et
$i : A \to B$, le \emph{produit-boîte} est
\[
j \boxtimes i : A' \otimes B \amalg_{A' \otimes A} B' \otimes A
\longrightarrow B' \otimes B .
\]
Un carré de $j$ vers $\alpha(i,f)$ est la même chose qu'un carré de
$j \boxtimes i$ vers $f$, et les diagonales se correspondent :
\[
j \pitchfork \alpha(i,f) \iff j \boxtimes i \pitchfork f, \qquad
\alpha(i', \alpha(i,f)) \simeq \alpha(i' \boxtimes i, f) .
\]
Il en tire le tableau de la page 11 :
\[
\begin{array}{rcl}
C \otimes T \subset C & \iff & \mathbf{Hom}(T, ?) \text{ préserve } TF, \\
TC \otimes T \subset TC & \iff & \mathbf{Hom}(T, ?) \text{ préserve } F, \\
C \boxtimes C \subset C & \iff & \bigl( i \in C,\ f \in TF \Rightarrow \alpha(i,f) \in TF \bigr), \\
C \boxtimes TC \subset TC & \iff & \bigl( i \in TC,\ f \in F \Rightarrow \alpha(i,f) \in TF \bigr), \\
TC \boxtimes C \subset TC & \iff & \bigl( i \in C,\ f \in F \Rightarrow \alpha(i,f) \in F \bigr).
\end{array}
\]
Chaque ligne se lit sur l'équivalence précédente : par exemple
$\alpha(i,f) \in TF = C_*$ pour tout $i \in TC$, $f \in F$ équivaut à
$j \boxtimes i \in F^* = TC$ pour tout $j \in C$, $i \in TC$.\footnote{Le
membre de gauche de la quatrième ligne est noyé sous des hachures, sous
lesquelles on lit « $TC \boxtimes TC$ » ou, avec doute, « $C \boxtimes TC$ ».
C'est $C \boxtimes TC$ que l'équivalence demande : la variable extérieure $j$
parcourt $C$, puisque $TF = C_*$. Dans le membre de droite, $TF$ est écrit
sur une autre lettre. Les deux premières lignes sont écrites sur
« stable by $? \otimes T$ », biffé.}

Les deux premières lignes sont des cas particuliers des trois autres, ce que la
page ne dit pas. Comme $\varnothing \otimes X = \varnothing$ dans un topos, on
a $(\varnothing \to T) \boxtimes i = i \otimes T$ et
$\alpha(\varnothing \to T, f) = \mathbf{Hom}(T, f)$ ; et
$\varnothing \to T$ est un monomorphisme.\footnote{La remarque est la nôtre.
Elle explique la mention marginale « under suitable assumptions » de la page
6 : la condition (iii) n'équivaut pas à la stabilité par $? \times T$ ; elle
équivaut à la stabilité des classes par produit-boîte (troisième à cinquième
lignes), dont la stabilité par $? \times T$ est le cas
$j = (\varnothing \to T)$.} La condition (iii), avec ses trois lignes, est ce
qu'on appelle aujourd'hui l'axiome du produit-boîte d'une catégorie de
modèles monoïdale, ou l'axiome SM7 de Quillen quand $\mathcal{E}$ est enrichie
sur elle-même.\footnote{Quillen, \emph{Homotopical Algebra} (1967), pour SM7 ;
la forme monoïdale est chez Hovey (1999) et Schwede--Shipley (2000).
Références citées de mémoire.}

\pagerange{11}{13}
\section*{L'exemple : cofibrations monomorphes et cylindres (pages 11 et 13)}

Soit $\mathcal{E}$ un topos, $\otimes$ le produit cartésien, $\mathbf{Hom}$ le
$\mathrm{Hom}$ interne, et $C$ la classe des monomorphismes.\footnote{Pour que
$C = (TF)^*$ avec $TF = (C_0)_*$, il faut que les monomorphismes soient la
saturation d'un petit ensemble $C_0$. C'est vrai dans tout topos de
Grothendieck (Cisinski, cité de mémoire) ; la page prend $C$ égal aux
monomorphismes sans le dire.} Pour un monomorphisme $i_0 : A_0 \to B_0$ et un
objet $I$ muni d'une section $\varepsilon : e \to I$, avec $I \to e$ dans $W$,
soit
\[
h(i_0, I, \varepsilon) : (B_0 \times e) \cup (A_0 \times I)
\subset B_0 \times I ,
\]
où $B_0 \times e$ désigne l'image de $\mathrm{id} \times \varepsilon$ ;
c'est le produit-boîte $i_0 \boxtimes \varepsilon$. On prend pour $(TC)_0$
l'ensemble de ces flèches, $I$ parcourant un ensemble donné (« a bunch »)
d'intervalles munis de leur section.\footnote{Il écrit $e$ pour la section
comme pour l'objet final ; on écrit $\varepsilon$. « $W$-aspheric » est son
mot pour « $I \to e$ est dans $W$ » ; « object » est lu avec doute. Le
« bunch » est en marge : « given bunch of $I$'s with given sections ».}

Les vérifications suivent le tableau.
\begin{itemize}
\item $C \times T \subset C$ est trivial, donc $\mathbf{Hom}(T, ?)$ préserve
$TF$.\footnote{La page écrit, avec doute, « $F$ stable » ; la lettre qui
précède $F$, s'il y en a une, est à peine marquée. C'est $TF$ que la première
ligne du tableau donne.}
\item $(TC)_0 \times T \subset (TC)_0$, car $h(i_0, I, \varepsilon) \times T =
h(i_0 \times \mathrm{id}_T, I, \varepsilon)$ ; donc $\mathbf{Hom}(T, ?)$
préserve $F = ((TC)_0)_*$, et $TC \times T \subset TC$.
\item $C \boxtimes C \subset C$ : le produit-boîte de deux monomorphismes d'un
topos est un monomorphisme, parce qu'une réunion de deux sous-objets y est la
somme amalgamée le long de leur intersection.
\item $TC \boxtimes C \subset TC$ et $C \boxtimes TC \subset TC$ : sur les
générateurs, par associativité du produit-boîte,
\[
i' \boxtimes h(i_0, I, \varepsilon) = h(i' \boxtimes i_0, I, \varepsilon),
\]
qui est encore dans $(TC)_0$ puisque $i' \boxtimes i_0$ est un
monomorphisme. La page le calcule sur deux carrés et conclut en
français : « On trouve … vérifier … et la formule symétrique, donc OK. »
\end{itemize}

Il reste à passer des générateurs à $TC$, et la page ne le fait
pas.\footnote{Les deux carrés de la page 13 portent les étiquettes
$h(i_0 \times \mathrm{id}_{A'}, I, e)$ et $h(i_0 \times \mathrm{id}_{B'}, I,
e)$, reliées aux flèches verticales par des accolades ; plusieurs lettres du
premier carré, page 11, sont surchargées. « Vérifier » et « OK » sont lus
avec doute.} Pour $i' \in C$ fixé, la classe des $j$ tels que
$j \boxtimes i' \in TC = F^*$ est celle des $j$ tels que
$j \pitchfork \alpha(i', f)$ pour tout $f \in F$ : c'est une classe de la
forme $\Psi^*$, donc saturée à gauche. Elle contient $(TC)_0$, donc
$((TC)_0)_*{}^* = TC$. Les trois dernières lignes du tableau sont donc
satisfaites, et la condition (iii) aussi.\footnote{Cet argument est le nôtre.
Le procédé (produit-boîte d'un monomorphisme par l'inclusion d'une extrémité
d'un intervalle, puis saturation) est celui qui engendre les « extensions
anodines » de Cisinski (\emph{Théories homotopiques dans les topos}, 2002 ;
\emph{Les préfaisceaux comme modèles des types d'homotopie}, Astérisque 308,
2006). Cisinski est l'élève de Maltsiniotis et ses travaux partent de
\emph{Pursuing Stacks} ; mais le manuscrit ne peut pas les connaître, et l'on
ne dit pas que l'un dérive de l'autre.}

\emph{Le programme} (page 13). « We do not know yet » M1, l'existence
des deux factorisations $f = pi$, ni M5, la saturation de $\widetilde W$ ; or
il suffirait de vérifier a) et b) pour avoir une catégorie de modèles
fermée.\footnote{Sa numérotation n'est pas celle de Quillen, chez qui, sauf
erreur de mémoire, M1 est le relèvement, M2 la factorisation et M5 le
deux-sur-trois ; à la page 24 il désigne par « (M1) (M2) » les deux à la
fois.} La factorisation est « the usual argument for constructing injectives », moyennant un petit ensemble de générateurs et des limites
inductives filtrantes. Il ne veut pas en écrire la forme générale pour le
moment et la suppose acquise ; elle viendra aux pages 33 à 39. Le but est
maintenant d'identifier $\widetilde W$, et même de prouver $\widetilde W = W$.
Il faut alors rendre explicite la dépendance de $(TC)_0$ en $W$. On voudrait
prendre pour « bunch » tous les intervalles $W$-asphériques, mais ce bunch
n'est pas petit, ce qui gênera la factorisation : on en prendra un plus petit
et l'on cherchera ce qu'il doit contenir.

\pagerange{15}{19}
\section*{Comparer $\widetilde W$ et $W$ (pages 15 à 19)}

\subsection*{L'inclusion $\widetilde W \subset W$, sous cinq axiomes}

Comme $\widetilde W = TF \circ TC$ et que $W$ est stable par composition, il
suffit de montrer $TF \subset W$ et $TC \subset W$.

\emph{Les fibrations triviales.} Puisque $TF = C_*$ ne dépend que des
monomorphismes, l'inclusion $TF \subset W$ peut être prise comme axiome sur
$W$. Appelons intervalle \emph{séparant} un objet $I$ muni de deux points
$e \rightrightarrows I$ dont l'intersection est l'objet initial. Si un tel $I$
est injectif, c'est-à-dire si $I \to e$ est dans $TF$, l'axiome équivaut à :
$I \to e$ est \emph{universellement} dans $W$ (« $\mathrm{U}W$ »), au sens où
$I \times X \to X$ est dans $W$ pour tout $X$. L'objet de Lawvere $L =
\Omega$, qui classe les sous-objets, est un intervalle séparant injectif, et
l'axiome devient :
\[
\text{(1)} \qquad (\Omega \to e) \in \mathrm{U}W .
\]
Les deux sens de l'équivalence se vérifient ainsi.\footnote{La page énonce
l'équivalence et les deux faits qui suivent ; la démonstration est la nôtre.
La formule « $(I \to e) \in \mathrm{U}W$, i.e. $(I \times X \to X)$ is in
$W$ for any $X$ » est encadrée puis biffée ; on garde la définition qu'elle
donne de $\mathrm{U}W$. « Really » (« Thus, the axiom is really ») et
« Lawvere element », au-dessus du $L$, sont lus avec doute. La note oblique de
la marge gauche, lue par fragments, finit par « but this is always so, because
the [Lawvere] object is injective !!! ».} Si $TF \subset W$, alors $I \times X
\to X$, changement de base de $I \to e \in TF$, est dans $TF \subset W$.
Inversement, supposons (1) (ou la même propriété pour un intervalle séparant
$I$) et soit $f : X \to Y$ dans $TF$. Le relèvement de $\varnothing \to X$
fournit une section $s$ de $f$. Comme les deux points de $I$ sont disjoints,
$X \sqcup X = X \times \partial I \to X \times I$ est un monomorphisme, et le
relèvement fournit une homotopie $H : X \times I \to X$ au-dessus de $Y$
entre $sf$ et $\mathrm{id}_X$. Les deux sections $X \to X \times I$ sont dans
$W$ par (1) et le deux-sur-trois, donc $H$ et $sf$ le sont ; et $fs =
\mathrm{id}_Y$, $sf \in W$ donnent $f \in W$ par la condition c) de la page
4. La page le résume ainsi : toute flèche de $TF$ est une $I$-homotopie,
et toute $L$-homotopie est dans $W$.

\emph{Les cofibrations triviales : les générateurs.} Il faut d'abord
$(TC)_0 \subset W$, et pour cela deux axiomes :
\[
\text{(2)} \qquad (I \to e) \in \mathrm{U}W \ \text{ pour tout } I \text{ du
bunch},
\]
ce qui équivaut, par le deux-sur-trois, à $A \times e \to A \times I$ dans $W$
pour tout $A$ ; et
\begin{enumerate}
\item[(3)] si $S \hookrightarrow T$ est un monomorphisme et
$g : S \hookrightarrow S'$ un monomorphisme dans $W$, la flèche
$T \to T' = S' \amalg_S T$ est dans $W$.
\end{enumerate}
Alors, dans le diagramme

\begin{tikzcd}[column sep=small, row sep=small, nodes={font=\scriptsize}]
A_0 \times e \arrow[r, hook] \arrow[d, "W"'] & B_0 \times e \arrow[d, hook] \arrow[dr, "W"] & \\
A_0 \times I \arrow[r, hook] & (B_0 \times e) \cup (A_0 \times I) \arrow[r, hook] & B_0 \times I
\end{tikzcd}

le carré est cocartésien, la flèche verticale médiane est dans $W$ par (3), la
diagonale par (2), et $h(i_0, I, \varepsilon)$ par le
deux-sur-trois.\footnote{La flèche verticale médiane porte sur la page
« $f$ » et « $W$ » ; le premier terme est précédé d'une parenthèse biffée. La
condition (3) est encadrée sur la page 17, où il la dit nécessaire de toute
façon pour $W = \widetilde W$.}

\emph{Les cofibrations triviales : la saturation.} Soit $f : A \to B$ dans
$TC$. On le factorise en $f = pi$ avec $p \in F$ et $i$ obtenu par l'argument
du petit objet, c'est-à-dire comme composé transfini de changements de cobase
de sommes de flèches de $(TC)_0$. Par (3) et l'axiome
\begin{enumerate}
\item[(4)] $W$ est stable par limites inductives filtrantes,
\end{enumerate}
$i$ est dans $W$.\footnote{La page dit « when we use factorization » ; il faut
que ce soit la factorisation cellulaire, sans quoi (3) et (4) ne s'appliquent pas
à $i$. Les sommes s'obtiennent par (3) pour un nombre fini de termes, puis par
(4). Le chiffre entouré après « Thus we need » est mal formé ; on le lit (4).}
Comme $f \pitchfork p$, il existe $s$ avec $sf = i$ et $ps = \mathrm{id}_B$,
et $f$ est un rétracte de $i$ :

\begin{tikzcd}
A \arrow[r, hook, "i"] \arrow[d, "="'] & B' \arrow[d, "p"] \\
A \arrow[r, hook, "f"'] & B \arrow[u, bend right=40, "s"']
\end{tikzcd}

Il faut donc enfin
\begin{enumerate}
\item[(5)] $W$ est stable par rétractes (« direct factors »),
\end{enumerate}
condition nécessaire pour $W = \widetilde W$, puisque $\widetilde W$ l'est
dès qu'on a une structure de modèles.\footnote{Il demande si (5) ne serait pas
conséquence de la saturation. Elle ne l'est pas des conditions a), b), c) de
la page 4 : voir la note de la page 26. Les deux figures de la marge gauche
portent, l'une, « $W$ » et « (and $TC$) » sur $i$, « $F$ » sur $p$ et
« $TC$ » sur $f$ ; l'autre dessine l'identité de gauche comme deux flèches
opposées entre deux signes « $=$ ».} « Thus (1) to (5) imply $\widetilde W \subset W$. »

\subsection*{L'inclusion $W \subset \widetilde W$, qui reste ouverte}

Pour $\widetilde W \subset W$ il a fallu restreindre les intervalles ((2)) ; pour
que $\widetilde W$ soit assez grande il faut que le bunch soit assez gros, et
la condition naturelle serait qu'il contienne un intervalle séparant. On en a
d'éligibles, $W$-asphériques et séparants, comme $\Omega$.

Soit $f \in W$ factorisée en $f = pi$, $i \in TC$, $p \in F$. On sait déjà que
$i \in W$, donc $p \in W$. Les trois inclusions
\[
W \subset \widetilde W, \qquad F \cap W \subset TF, \qquad C \cap W \subset TC
\]
sont donc équivalentes, moyennant les deux factorisations.\footnote{La page
dit la deuxième et la troisième équivalentes, et équivalentes à la première,
dans une phrase en partie illisible. La vérification est immédiate : si
$F \cap W \subset TF$, la factorisation ci-dessus met $f$ dans
$\widetilde W$ ; si $C \cap W \subset TC$, on factorise $f = pi$ avec
$i \in C$, $p \in TF \subset W$, et $i \in C \cap W \subset TC$ ; enfin
$W \subset \widetilde W$ donne $F \cap W \subset F \cap \widetilde W = TF$ par
le lemme de la page 6.} La troisième dit que toute flèche ayant la propriété
de relèvement à droite par rapport aux $h(i_0, I, \varepsilon)$ l'a par
rapport à \emph{tous} les monomorphismes qui sont dans $W$. La page
s'interrompt sur « [ ? ??? ] ». Il voudrait l'avoir au moins dans deux cas.\footnote{Il souligne $W$ pour le localisateur fondamental ; on écrit
$\mathcal{W}$. Il écrit $T$ pour la catégorie test de b), lettre déjà prise
par l'objet de la page 6 ; on écrit $B$. « Locale » est lu avec doute. Le cas
a) est aujourd'hui un théorème de Cisinski (Astérisque 308, 2006) : pour un
localisateur fondamental accessible $\mathcal{W}$ et une catégorie
$\mathcal{W}$-test locale $A$, $\mathcal{W}_A$ est la classe des équivalences
faibles d'une structure de modèles sur $\widehat{A}$ dont les cofibrations sont
les monomorphismes. L'énoncé est cité de mémoire, et l'on ne dit pas sous
quelle forme le cas b) y figure. La note oblique de la marge gauche (« but
this cannot possibly be enough, as it would … given just one particular … »)
ne se lit que par fragments.}
\begin{enumerate}
\item[a)] $\mathcal{E} = \widehat{A}$ et $W = \mathcal{W}_A$, où
$\mathcal{W}$ est un localisateur fondamental et $A$ une catégorie
$\mathcal{W}$-test (locale ?).
\item[b)] $A = A_0 \times B$ avec $B$ une catégorie test, de sorte que
$\widehat{A} \simeq \mathbf{Hom}(A_0^{\mathrm{op}}, \widehat{B})$, et
$W = \mathcal{W}_{A/A_0}$, défini fibre à fibre par $\mathcal{W}_B$.
\end{enumerate}

\pagerange{21}{21}
\section*{Une présentation alternative (page 21)}

La page 21 (sa page 9) reprend tout depuis l'autre bout : « Alternative
presentation of theory, with possibly slightly less strong result ». On se
donne cette fois $W$, on veut une structure de modèles dont les cofibrations
sont les monomorphismes et les équivalences faibles les flèches de $W$, et
l'on pose directement
\[
C = \text{mono}, \quad TC = C \cap W, \quad F = (TC)_*, \quad TF = C_* .
\]
On sait toujours que $(TF)^* = C$. Restent quatre problèmes.\footnote{« Kan- » suivi d'un mot illisible désigne, deux fois, les flèches de
$TF$. La page dit que c) découle « de 3) et 4) », et (a) « de 5) » ; la
transcription note que ces conditions ne sont pas posées dans le lot, mais ce
sont les axiomes (3), (4), (5) des pages 15 à 17. La condition (d), « quite
important for me », est dite duale d'une condition lue avec doute « (b) ».
C'est ce qu'on appelle aujourd'hui la propreté à droite, duale de la
stabilité (3) par changement de cobase.}
\begin{enumerate}
\item[(a)] $F^* = TC$. Comme $F^* \subset (TF)^* = C$, cela revient à : tout
monomorphisme ayant la propriété de relèvement à gauche par rapport à $F$ est
dans $W$. C'est une propriété de saturation de $W$ relativement aux
monomorphismes.
\item[(b)] $TF = W \cap F$. Les hypothèses déjà faites ((1), page 15) donnent
$TF \subset W$, donc $TF \subset W \cap F$ ; il reste $W \cap F \subset TF$.
\item[c)] Les factorisations, la seconde avec $i \in C \cap W$ ; elles
« shouldn't be too hard » par des arguments standard, et il
les admet.
\item[(d)] Si $f : X \to Y$ est dans $F$ et $Y' \to Y$ dans $W$, alors
$X' = X \times_Y Y' \to X$ est dans $W$.
\end{enumerate}

Pour (a), l'argument de rétracte de la page 17 vaut tel quel. Si $f$ est un
monomorphisme de $F^*$, on le factorise en $f = pi$ avec $i \in W \cap C$,
$p \in F$. Le relèvement donne $s$ avec $sf = i$, $ps = \mathrm{id}$, donc $f$
est un rétracte de $i$, et $f \in W$ par (5).\footnote{La lettre qui suit « Thus
for » est tracée comme le c) qui précède ; l'argument est celui qui répond à
(a). Le petit diagramme est dans la marge inférieure.}

\begin{tikzcd}
& X' \arrow[d, "p \in F"] \\
X \arrow[ur, "i \in W \cap C"] \arrow[r, "f"'] & Y \arrow[u, bend right=40, "s"']
\end{tikzcd}

La différence avec les pages 4 à 19 est de principe : là, $TC$ était engendrée
par $(TC)_0$ et il fallait comparer $\widetilde W$ à $W$ ; ici $TC$ est
définie par $W$, et la question de la génération reviendra à la page 31.

\pagerange{24}{27}
\section*{Classes saturées, factorisation, structures de modèles (pages 24 à 27)}

\subsection*{Saturations et systèmes de factorisation faibles}

Soit $M$ une catégorie. Une classe $\Phi \subset \mathrm{Fl}(M)$ est
\emph{saturée à gauche} si $\Phi = \Psi^*$ pour une classe $\Psi$,
c'est-à-dire si $\Phi = (\Phi_*)^*$ ; une classe $\Psi$ est \emph{saturée à
droite} si $\Psi = (\Psi^*)_*$. On a toujours $\Phi \subset (\Phi_*)^*$ et
$\Psi \subset (\Psi^*)_*$, et $\Phi \mapsto \Phi_*$, $\Psi \mapsto \Psi^*$
sont des bijections décroissantes, inverses l'une de l'autre, entre classes
saturées à gauche et classes saturées à droite. Une classe saturée à droite
est stable par composition, changement de base et rétractes ; dualement à
gauche.\footnote{La dernière phrase est une note oblique de la marge gauche,
où « right », « base ch. » et « retracts » sont lus avec doute. Il dessine son
$\Phi$ tantôt rond, tantôt en une forme cursive proche d'un L de ronde, dans
les mêmes formules.}

Deux classes $\Phi, \Psi$ sont \emph{orthogonales} si $\Phi \subset \Psi^*$,
ou, ce qui revient au même, $\Psi \subset \Phi_*$. La paire a la
\emph{propriété de factorisation} si toute flèche $f$ s'écrit $f = pi$ avec
$i \in \Phi$, $p \in \Psi$.

\textbf{Proposition.} Soit $(\Phi, \Psi)$ une paire orthogonale ayant la
propriété de factorisation. Alors $\Phi = \Psi^*$ si et seulement si $\Phi$
est stable par rétractes, et $\Psi = \Phi_*$ si et seulement si $\Psi$ l'est.
En effet, si $h \in \Psi^*$, on factorise $h = pi$, et $h \pitchfork p$ fait
de $h$ un rétracte de $i \in \Phi$, comme à la page 17.

Il appelle « paire de Quillen » une paire orthogonale ayant la propriété de
factorisation dont les deux classes sont stables par rétractes : c'est ce
qu'on appelle aujourd'hui un \emph{système de factorisation
faible}.\footnote{Le nom « paire de Quillen » désigne aujourd'hui tout autre
chose : une adjonction entre catégories de modèles qui préserve cofibrations
et cofibrations triviales. Sur la page, c'est la paire $(C, TF)$ ou $(TC, F)$
d'une structure de modèles. On ne garde pas son nom pour éviter la confusion.
Il dit aussi « dual pair of subsets (localizers) » pour une paire
$(\Psi^*, \Psi)$.} Il appelle « triple de Quillen » une donnée $(C, F, W)$
telle que :
\begin{enumerate}
\item[a)] $W$ est faiblement saturée et stable par rétractes (M5) ;
\item[b)] $(C, TF)$ et $(TC, F)$ sont des systèmes de factorisation faibles,
où $TF = F \cap W$ et $TC = C \cap W$ (M1, M2).
\end{enumerate}
C'est une \emph{structure de modèles} au sens de Quillen ; il en résulte que
$f \in W$ si et seulement si $f = pi$ avec $i \in TC$, $p \in TF$, et, si $M$
a des limites projectives et inductives finies (M0), que $W$ est fortement
saturée.\footnote{L'insertion « it would not be necessary here to assume the
direct factor condition on $W$, if we assume b) » est juste : la stabilité de
$W$ par rétractes découle des autres axiomes (Joyal et Tierney, cité de
mémoire). La forte saturation est un résultat de Quillen ; la condition M0 est
en marge, où « limits » et « finite inverse limits » sont lus avec doute.}

\pagerange{26}{26}
\subsection*{Un brouillon : la saturation (page 26)}

Le feuillet n'a pas de numéro de sa main et il est barré de deux longues
diagonales au crayon. Il commence par les complexes de Kan : si $Y$ est de
Kan, $\mathbf{Hom}(I, Y) \to Y \times Y$ est une fibration de Kan ; est-ce une
fibration de Serre ? La question est aussitôt posée fibre à fibre, au-dessus
d'un $a \to Y \times Y$.\footnote{La première assertion est juste ; c'est un
cas de l'axiome SM7 pour les ensembles simpliciaux. La seconde question n'est
pas développée ; les notations $\mathbf{Hom}_{\alpha, \beta}(I_a, Y_a)$ sont
celles de la page.}

Puis trois conditions sur une classe $W$ : a) $\mathrm{Iso}_M \subset W$ ; b)
le deux-sur-trois ; c) $W$ stable par rétractes. Elles donnent la saturation
faible. Donnent-elles la saturation, $fg, gf \in W \Rightarrow f, g \in W$ ?
Il compare
\begin{enumerate}
\item[(i)] $gf = \mathrm{id}$, $fg \in W$ $\Rightarrow$ $f, g \in W$ ;
\item[(ii)] $gf \in W$, $fg \in W$ $\Rightarrow$ $f, g \in W$ ;
\item[(iii)] $W$ stable par rétractes ;
\end{enumerate}
et note (ii) $\Rightarrow$ (i) $\Leftarrow$ (iii). La seconde implication
vient de ce que, si $gf = \mathrm{id}_A$, $f$ est un rétracte de $fg$ :

\begin{tikzcd}
A \arrow[r, "f"] \arrow[d, "f"'] & B \arrow[r, "g"] \arrow[d, "fg"] & A \arrow[d, "f"] \\
B \arrow[r, "\mathrm{id}"'] & B \arrow[r, "\mathrm{id}"'] & B
\end{tikzcd}

et $g$ est alors dans $W$ par le deux-sur-trois. La question elle-même reste
ouverte sur la page, et la réponse est non.\footnote{Le contre-exemple est le
nôtre. Soit la catégorie libre engendrée par deux flèches $f : A \to B$,
$g : B \to A$ ; ses flèches sont les mots alternés, de longueur $\ell$, avec
$\ell(vu) = \ell(u) + \ell(v)$. Soit $W$ la classe des mots de longueur paire.
Elle contient les identités, seuls isomorphismes ; elle a le deux-sur-trois par
additivité de la longueur ; elle est stable par rétractes, car seules les
identités ont une section dans une catégorie libre. Elle a donc a), b), c), et
pourtant $gf, fg \in W$ alors que $f \notin W$. La propriété (ii) est un cas de
ce qu'on appelle aujourd'hui le deux-sur-six (Dwyer, Hirschhorn, Kan, Smith,
2004), qui vaut pour toute classe fortement saturée. Autour de ces lignes, des
esquisses : un encadré « a) b) c$_1$) (c$_2$) (c$_3$) », et des homotopies
$gf \sim_h \mathrm{id}$, $fg \sim_h \mathrm{id}$ qui mettent $gf$ et $fg$ dans
$W$ ; un dernier schéma n'est pas restitué.}

\pagerange{27}{27}
\subsection*{Conditions équivalentes (page 27, sa page « 10 bis »)}

Soit $(C, F, W)$ avec $TF = F \cap W$ et $TC = C \cap W$, et posons
$\widetilde{TF} = C_*$, $\widetilde{TC} = F^*$. On suppose
\[
\text{(a)} \ \ C = (TF)^*, \ F = (TC)_*, \qquad
\text{(b)} \ \ W = \{ pi \mid p \in \widetilde{TF},\ i \in \widetilde{TC} \}.
\]
C'est la construction $\widetilde W$ de la page 6, faite avec les paires
$(F^*, F)$ et $(C, C_*)$. Comme $C$ et $F$ sont saturées, $\widetilde{TF}
\subset F$ et $\widetilde{TC} \subset C$ ; et (a) donne $TF \subset
\widetilde{TF}$, $TC \subset \widetilde{TC}$. Il suffit donc que
$\widetilde{TF} \subset W$ et $\widetilde{TC} \subset W$, ce qui équivaut à
$\widetilde{TF} = TF$ et $\widetilde{TC} = TC$, c'est-à-dire à ce que
$(C, TF)$ et $(TC, F)$ soient des systèmes de factorisation faibles. D'où les
conditions équivalentes : (i) (a) et (b), une catégorie de modèles fermée au sens
de Quillen ; (ii) $(C, TF)$ et $(TC, F)$ sont des systèmes de factorisation
faibles ; et il faut encore (iii) $C$, $F$, $W$ stables par
rétractes.\footnote{La factorisation de toutes les flèches, admise depuis la
page 21, est implicite dans (i). La place de (iii) dans l'équivalence n'est
pas claire sur la page ; il ajoute aussitôt : « I've to check » que (i) ou (ii)
entraîne que $W$ est stable par rétractes. Le listing sur lequel il écrit
porte imprimés « 06/09/82-09:14:22 » et « CFT 1.09 (05/10/82) ».}

Quillen démontre plus : $W$ est fortement saturée, du moins quand $M$ a des
limites inductives et projectives finies. La page le ramène aux objets
fibrants et cofibrants. Par des factorisations successives, $f : X \to Y$ est
relié à une flèche $g : X_2 \to Y_2$ entre objets fibrants et cofibrants, de
sorte que $\gamma(f)$ est un isomorphisme si et seulement si $\gamma(g)$ en
est un, et $f \in W$ si et seulement si $g \in W$. Donc $W$ est fortement
saturée si et seulement si, pour $f$ entre objets fibrants et cofibrants et
$f \in C$, $\gamma(f)$ isomorphisme entraîne $f \in W$ ; et de même avec $f \in
F$.\footnote{Le diagramme de remplacement de la page passe par l'objet
initial, tracé $\varphi$, et porte deux flèches en pointillé marquées
« $\varepsilon$ » et « $F$ » qu'on ne restitue pas. C'est le théorème de
Whitehead des catégories de modèles : une flèche entre objets fibrants et
cofibrants inversée dans la catégorie homotopique est une équivalence
d'homotopie, donc une équivalence faible.}

\pagerange{29}{31}
\section*{Construire une structure de modèles à partir de $(C, W)$ (pages 29 et 31)}

La page 29 (sa page 11) donne la recette. On part d'une classe
$C \subset \mathrm{Fl}(M)$ stable par rétractes et l'on vérifie la
factorisation pour la paire $(C, C_*)$, qui est alors un système de
factorisation faible ; on pose $TF = C_*$. On prend ensuite une classe
$W \supset TF$, faiblement saturée et stable par rétractes. Alors $TC = C
\cap W$ est stable par rétractes ; on pose $F = (TC)_*$, qui contient $TF$,
et l'on vérifie la factorisation pour $(TC, F)$. Il reste $TF = F \cap W$. Une
inclusion est acquise. Pour l'autre, soit $f : X \to Y$ dans $F \cap W$,
factorisé en $f = pi$ avec $i \in C$, $p \in TF$. Comme $f, p \in W$, on a
$i \in W$, donc $i \in TC$. Le relèvement de $i$ par rapport à $f$ donne
$s : X' \to X$ avec $si = \mathrm{id}_X$ et $fs = p$, de sorte que $f$ est un
rétracte de $p$, et $f \in TF$.

\begin{tikzcd}
X \arrow[r, "="] \arrow[d, "i"'] & X \arrow[d, "f \in F"] \\
X' \arrow[r, "p"'] \arrow[ur, dashed, "s"] & Y
\end{tikzcd}

\textbf{Proposition.} Pour qu'une paire $(C, W)$ provienne d'une structure de
modèles $(C, F, W)$, il faut et il suffit que :\footnote{Que $C = (C_*)^*$ et $C \cap W = ((C \cap W)_*)^*$ découle de 1) et
2) par la proposition de la page 24 ; la page ne le dit pas. La
caractérisation est aujourd'hui standard : une structure de modèles est
déterminée par ses cofibrations et ses équivalences faibles. Il dit « mildly
saturated ».}
\begin{enumerate}
\item[1)] $C$ et $W$ soient stables par rétractes, et $W$ faiblement
saturée ;
\item[2)] les paires $(C, C_*)$ et $(C \cap W, (C \cap W)_*)$ aient la
propriété de factorisation ;
\item[3)] $W \supset C_*$.
\end{enumerate}

« In the case of interest to us », 1) est gratuit, 3) presque (il vient de
l'hypothèse « test »), et 2) pose un petit problème technique. Il s'intéresse
aussi, au-delà des factorisations, à la propriété qui caractérise un triple
de Quillen \emph{strict} :
\begin{enumerate}
\item[4)] $W$ est stable par changement de base le long des $g \in F$ et par
changement de cobase le long des $g \in C$.
\end{enumerate}
C'est ce qu'on appelle aujourd'hui une structure de modèles
\emph{propre}.\footnote{Le terme est de Bousfield et Friedlander (1978), cité
de mémoire. La condition de base est la (d) de la page 21, celle de cobase
la (3) de la page 15, sous une forme où la flèche le long de laquelle on
change de cobase est une cofibration quelconque.}

\emph{NB} (page 31). C'est la réflexion de méthode du dossier. De son point de vue, $W$ est donnée d'avance, de façon assez
tangible, par un critère homologique, très maniable. Chez Kan et Quillen, $W$
était plus cachée : on y arrive par la paire $(C, TF)$ et les fibrations de
Kan $F \supset TF$, une construction quelque peu ad hoc à partir d'un ensemble
générateur $(TC)_0$, qui exprime en gros la propriété de relèvement de Kan.
Reste à comprendre, pour une catégorie test générale, le rapport entre cette
construction et la description directe $TC = C \cap W$, c'est-à-dire pourquoi
$C \cap W$ est exactement la \emph{saturation} de $(TC)_0$. Peut-être faut-il
regarder le traitement par Gabriel et Zisman des « extensions
anodines ».\footnote{« ss structures » et « a-priori » sont lus avec doute.
Gabriel et Zisman, \emph{Calculus of Fractions and Homotopy Theory} (1967),
montrent que, pour les ensembles simpliciaux, les monomorphismes qui sont des
équivalences faibles sont exactement la saturation des inclusions de cornets.
Pour une catégorie test, la question est celle à laquelle répond la théorie de
Cisinski (2006), où les cofibrations triviales sont la saturation d'un petit
ensemble d'extensions anodines (cité de mémoire).}

\pagerange{33}{39}
\section*{Factorisation (pages 33 à 39)}

Le titre est le sien, souligné, en oblique dans la marge supérieure de la page
33 (sa page 13). Ces pages fournissent ce qui avait été admis aux pages 13 et
21. C'est aussi là que se décide la question de la page 15 : le bunch
d'intervalles doit être petit, parce que $(TC)_0$ doit l'être.

\subsection*{La construction}

Soit $M$ une catégorie, $\Phi_0 \subset \mathrm{Fl}(M)$, $\Psi = (\Phi_0)_*$,
et $\Phi_0 \subset \Phi \subset \overline{\Phi}_0 = \Psi^*$. On cherche des
conditions sur $\Phi_0$ et $\Phi$ pour que toute flèche se factorise en
$f = pi$ avec $i \in \Phi$, $p \in \Psi$.\footnote{En interligne : « NB. we
don't use $\Phi \subset \overline{\Phi}_0$ » ; c'est juste, cette inclusion ne
sert qu'à l'orthogonalité. Un astérisque au-dessus de $\Phi_0$ renvoie,
semble-t-il, à la note marginale de la page 39 ; une note oblique sous le titre
renvoie à « p.~16 », sa page 16, qui est la page 39.}

Pour $f : X \to Y$ fixée, on considère tous les carrés commutatifs $(D)$

\begin{tikzcd}
A \arrow[r, "u"] \arrow[d, "g"'] & X \arrow[d, "f"] \\
B \arrow[r, "v"'] & Y
\end{tikzcd}

avec $g \in \Phi_0$, et l'on pose $X_D = X \amalg_A B$, puis $X_1(f)$ la somme
amalgamée sous $X$ de tous les $X_D$. On obtient une factorisation
$f = \Sigma_1(f) \circ i_1(f)$, avec $i_1(f) : X \to X_1(f)$ et
$\Sigma_1(f) : X_1(f) \to Y$, et $i_1(f) \in \Phi$ pourvu que :\footnote{Les insertions en bout des lignes b) et c) sont très serrées, et
leur place exacte est incertaine. Le carré supérieur de la construction de
$X_D$ est marqué « coc. ».}
\begin{enumerate}
\item[a)] $\Phi_0$ soit petit ;
\item[b)] les sommes amalgamées indexées par des ensembles de cardinal
$\leqslant \mathrm{card}\, \Phi_0 \times c$ existent dans $M$, où $c$ est le
cardinal des $\mathrm{Hom}(A, X) \times \mathrm{Hom}(B, Y)$ ;
\item[c)] $\Phi$ soit stable par composition, par changement de cobase et par
les limites inductives filtrantes indexées par des ensembles de cardinal
$\leqslant c$.
\end{enumerate}

Puis, par récurrence transfinie sur les ordinaux $\alpha < \alpha_0$,
\[
\Sigma_{\alpha+1}(f) = \Sigma_1(\Sigma_\alpha(f)), \qquad
i_{\alpha+1}(f) = i_1(\Sigma_\alpha(f)) \circ i_\alpha(f),
\]
et
\[
\Sigma_\alpha(f) = \varinjlim_{\alpha' < \alpha} \Sigma_{\alpha'}(f)
\quad \text{si } \alpha \text{ est un ordinal limite},
\]
$X_\alpha(f)$ étant la source de $\Sigma_\alpha(f)$. On suppose $\Phi$ stable
par les limites inductives qui interviennent aux ordinaux limites.

\emph{Le choix de $\alpha_0$.} Soit $\pi$ un cardinal tel que les sources et
buts des $g \in \Phi_0$ soient $\pi$-accessibles : $\mathrm{Hom}(A, -)$
commute aux limites inductives filtrantes « large with resp.\ to $\pi$ ». On
peut prendre $\pi = \sup_{g \in \Phi_0} \pi_g$. Si ces objets sont de
présentation finie, on arrête la récurrence au premier ordinal infini ; en
général on prend pour $\alpha_0$ le premier ordinal dont l'ensemble des
prédécesseurs est large par rapport à $\pi$.\footnote{Dans le cas de
présentation finie, la page écrit « $\pi = 0$ » d'un rond épais ; c'est
$\aleph_0$, la récurrence s'arrêtant à $\omega$. En général, on peut prendre
pour $\alpha_0$ un cardinal régulier $\kappa \geqslant \pi$ : l'ensemble des
ordinaux $< \kappa$ est $\kappa$-filtrant. Cette précision est la nôtre.} On
pose $\mathbf{X}(f) = X_{\alpha_0}(f)$, $i(f) = i_{\alpha_0}(f)$ et
$p(f) = \Sigma_{\alpha_0}(f)$.

Alors $i(f) \in \Phi$ et $p(f) \in \Psi$, sous l'hypothèse d) : les sources et
buts des $g \in \Phi_0$ sont accessibles, et, pour un cardinal $\pi$ assez
grand, tout système inductif filtrant de cardinal $\leqslant \pi$ à flèches de
transition dans $\Phi$ a une limite, les flèches $Z_i \to Z$ étant dans
$\Phi$. La raison est la suivante.\footnote{La page énonce la conclusion sans
cet argument, qui est le nôtre. Elle marque d'un double trait en marge la
fonctorialité qui suit.} Un carré de $g : A \to B$, $g \in \Phi_0$, vers
$p(f)$ a sa flèche $A \to X_{\alpha_0}(f) = \varinjlim X_\alpha(f)$ qui se
factorise par un $X_\alpha(f)$, $\alpha < \alpha_0$, par accessibilité de $A$.
Ce carré figure donc parmi ceux qui construisent $X_{\alpha+1}(f)$, et celui-ci
fournit la diagonale $B \to X_{\alpha+1}(f) \to X_{\alpha_0}(f)$. De plus la
factorisation est \emph{fonctorielle} en $f$.

\pagerange{37}{37}
\subsection*{L'énoncé (page 37)}

Supposons :
\begin{enumerate}
\item[a)] $\Phi_0$ petit, et les sources et buts de ses flèches accessibles ;
\item[b)] $M$ stable par limites inductives filtrantes ;
\item[(c)] pour toute $u : A \to X$ et toute $\alpha : A \to B$ dans $\Phi_0$,
la somme amalgamée $X' = X \amalg_A B$ existe et $\alpha' : X \to X'$ est
dans $\Phi$ ; et c') $\Phi$ est stable par composition ;
\item[(d)] $\Phi$ est stable par limites inductives filtrantes, au sens de la
page 35.
\end{enumerate}
Alors $(\Phi, \Psi)$ a la propriété de factorisation, et il existe même un
foncteur $\Sigma : \mathrm{Fl}(M) \to \mathrm{Fl}(M)$ au-dessus de $M$ (pour le
foncteur but $t$) et un morphisme fonctoriel $\mathrm{id} \to \Sigma$ dans
$[\mathrm{Cat}]_{/M}$, qui envoie $f : X \to Y$ sur $p(f) :
\mathbf{X}(f) \to Y$ par $(i(f), \mathrm{id}_Y)$, avec $i(f) \in \Phi$ et
$p(f) \in \Psi$.\footnote{Il souligne $\mathrm{Fl}$ dans
$\mathrm{Fl}(M)$, la catégorie des flèches ; et il précise en
interligne que ce sont des catégories « possibly “large” » au-dessus de
$M$. En marge de (d) : « It is enough to know it w.r.t.\ ordered sets $I$ where
sup of any two elements exists ».}

\begin{tikzcd}
\mathrm{Fl}(M) \arrow[rr, "\Sigma"] \arrow[dr, "t"'] & & \mathrm{Fl}(M) \arrow[dl, "t"] \\
& M &
\end{tikzcd}

\emph{Exemple.} Si $M$ a des sommes amalgamées et $\Phi = \overline{\Phi}_0 =
\Psi^*$, la condition (c) est triviale, une classe saturée à gauche étant
stable par changement de cobase. Pour (d), il propose une récurrence de Zorn sur
les paires $(J \subset I, u_J : \varinjlim_J A_j \to X)$ qui prolongent le
relèvement donné sur $A_{i_0}$, la surjectivité cherchée étant celle de
$\mathrm{Hom}(B, X) \to \mathrm{Hom}(g, f) = \varprojlim
\mathrm{Hom}(g_i, f)$.\footnote{L'argument de Zorn vaut pour un système
\emph{bien ordonné}, où chaque étape n'a qu'un prédécesseur immédiat ou est
une limite. Il ne passe pas tel quel à un système filtrant quelconque : deux
relèvements déjà choisis sur deux $A_j$ n'ont pas de raison de se recoller sur
un $A_i$ qui les majore. C'est bien la restriction que fait le lemme de la page
39, où la version filtrante est biffée au profit des systèmes ordinaux ; et
c'est tout ce que la construction des pages 33 à 35 utilise. « OK », « for
compositions », « Zorn », « induction » et « obvious » sont lus avec doute.
Le schéma qui pose le problème est en marge, avec la flèche cherchée
$\varphi : B \to X$ en pointillé.}

\pagerange{39}{39}
\subsection*{Le lemme et le théorème (page 39)}

\textbf{Lemme.} Soit $\Phi = \Psi^*$ une classe saturée à gauche. Alors :
\begin{enumerate}
\item[a)] $\Phi$ contient les isomorphismes ;
\item[b)] $\Phi$ est stable par composition et changement de cobase ;
\item[c)] $\Phi$ est stable par sommes amalgamées quelconques ;
\item[d)] $\Phi$ est stable par systèmes inductifs ordinaux : si $I$ est bien
ordonné et $(A_i)_{i \in I}$ un système inductif à flèches de transition dans
$\Phi$, continu aux ordinaux limites ($A_{i_0} = \varinjlim_{j < i_0} A_j$),
alors $A_{i_0} \to \varinjlim_I A_i$ est dans $\Phi$ pour tout $i_0$.
\end{enumerate}
En marge : a), b), d) entraînent c).\footnote{« Including » et le mot qui suit,
à la fin de b), sont lus avec doute et illisibles. Il dit « left
Q-saturated », Q pour Quillen.}

\textbf{Théorème.} Soit $M$ une catégorie stable par petites limites
inductives et accessible.\footnote{« Weakly » devant « accessible » est lu avec doute ; seule sert
l'accessibilité des sources et buts des flèches de $\Phi_0$. La note
marginale oblique, en regard de (1) et (2), à laquelle renvoient les astérisques,
commence par « We have to assume $\Phi_0$ small » ; le reste ne se lit que par
fragments, dont une variante a), b$'$), c), d) avec b$'$) la stabilité par
changement de cobase sans composition. L'énoncé est l'argument du petit objet
de Quillen (1967), et ce qu'on appelle aujourd'hui une classe à engendrement
cofibrant ; une catégorie stable par petites limites inductives et accessible,
avec un petit ensemble de générateurs, est ce que Gabriel et Ulmer (1971)
appellent localement présentable.}
\begin{enumerate}
\item[(1)] Soit $\Phi_0 \subset \mathrm{Fl}(M)$ petit, et
$\Phi_0 \subset \Phi \subset \overline{\Phi}_0 = \Psi^*$, $\Psi = (\Phi_0)_*$.
Si $\Phi$ satisfait a), b), c), d), la paire $(\Phi, \Psi)$ a la propriété de
factorisation ; si de plus $\Phi$ est stable par rétractes,
$\Phi = \overline{\Phi}_0$.
\item[(2)] Pour que $\Phi = \overline{\Phi}_0$, il faut et il suffit que $\Phi$
satisfasse a), b), c), d) et e) : la stabilité par rétractes. La
factorisation a lieu alors pour le système $(\Phi, \Psi)$.
\end{enumerate}

\emph{NB.} Dans (2), on supposait au préalable l'existence d'un $\Phi_0$
\emph{petit} ayant même saturation à gauche que $\Phi$. « Next step is to try
to get rid of this. » La suite n'est pas dans le dossier.\footnote{Sans
hypothèse, on ne peut pas s'en défaire : la question devient celle de savoir
quand une classe $C \cap W$ est engendrée par un petit ensemble, à laquelle
répondent, sous des hypothèses d'accessibilité sur $W$, le théorème de
J.~Smith (exposé par T.~Beke en 2000) et, pour les topos, Cisinski. Références
citées de mémoire.}

\pagerange{42}{45}
\section*{Les axiomes W7 et W8, et l'axiome « Der 6 » (pages 42 à 52)}

La dernière suite est en français et paginée de sa main de 6 à 18 ; ses pages 1
à 5 ne sont pas dans le dossier. Elle s'ouvre sur la fin, biffée, d'un
argument commencé avant : un carré de foncteurs $U \to \overline{U}$,
$U' \to \overline{U}'$ dont les deux flèches horizontales sont dans $W$, donc
« $q \in W$ ssi $q_U \in W$ cqfd ». Elle demande alors : « \emph{Dém.} de W7,
W 7 bis si $W = W_{\mathbb D}$ ? »\footnote{Le chiffre devant « bis » est
surchargé. La liste des axiomes W1 à W8, ni les énoncés de W7 et W8, ne sont
dans le dossier ; leur forme se lit sur l'argument, et sur les formes « bis »
des pages 51 et 52. De même les axiomes « Der 1 -- Der 5, 5' » de la page 48
sont cités sans être énoncés. La numérotation Der 1 à Der 5 est celle que la
théorie des dérivateurs a gardée, mais on ne prétend pas que la liste
qu'invoque la page coïncide avec la liste reçue.}

\emph{Notations de cette partie.} Un \emph{ouvert} d'une catégorie $X$ est un
crible : une sous-catégorie pleine $U$ telle que $x \to y$ et $y \in U$
entraînent $x \in U$.\footnote{C'est le sens que force la page, où
$U = \varphi^{-1}(\{0\} \times \Delta^1) = X_{/a}$.} On note
$Q = \Delta^1 \times \Delta^1$, et $\Psi = Q \smallsetminus \{(1,1)\}$,
l'ensemble ordonné $\{0 < a,\ 0 < b\}$ où $0 = (0,0)$, $a = (0,1)$, $b = (1,0)$ ;
$k : \Psi \subset Q$ est l'inclusion, qui ajoute à $\Psi$ l'élément final
$(1,1)$. On note $\Phi = \Psi^{\mathrm{op}} = (a \to 0 \leftarrow b)$. Pour
$\varphi : X \to I$ et $s \in I$, $X_{/s} = X \times_I I_{/s}$ ; la
construction de Grothendieck du foncteur $s \mapsto X_{/s}$ est notée
$\mathcal{G}_I(X)$.\footnote{Il la note d'un $\Psi$ épaissi affecté d'un
indice, « $\Psi_{Q}(X)$ », qui se confond avec l'ensemble ordonné
$\Psi$ ; on écrit $\mathcal{G}$.} $\mathbb D$ est un dérivateur,
$\mathcal{A} = \mathbb D(e)$ ; pour une catégorie $X$, $p_X : X \to e$, et
$\xi \in \mathcal{A}$,
\[
H^\bullet_{\mathbb D}(X, \xi) = p_{X*}\, p_X^*\, \xi \in \mathcal{A}
\]
est la cohomologie de $X$ à coefficients constants $\xi$, et $W_{\mathbb D}$
la classe des foncteurs $u : X \to Y$ tels que $u^* :
H^\bullet_{\mathbb D}(Y, \xi) \to H^\bullet_{\mathbb D}(X, \xi)$ soit un
isomorphisme pour tout $\xi \in \mathcal{A}$.\footnote{Cette définition n'est
pas sur ces pages ; elle se lit sur l'usage qu'en fait la page 45 (« ceci
étant vrai pour tout $\xi \in \mathcal{A}$, cela donne $i_0 \in W_{\mathbb D}
\Rightarrow i \in W_{\mathbb D}$ »). Il souligne $\mathbb D$ en un endroit, et
$\Psi$, $\Phi$, $Q$, $U_0$ partout.} Les objets de $\mathbb D(I)$ se lisent,
comme le dit la page 45, comme des diagrammes de type
$I^{\mathrm{op}}$.\footnote{C'est la convention des préfaisceaux, à laquelle
toutes les formules de ces pages obéissent : la valeur en $(1,1)$ d'un
$k_*(\gamma_0)$ y est une limite projective sur $\Psi^{\mathrm{op}}$, et
$H^\bullet_{\mathbb D}(Q, \gamma) \simeq \gamma(1,1)$ parce que $(1,1)$ est
final dans $Q$. Une bonne part de la littérature actuelle sur les dérivateurs
prend la convention opposée.}

\subsection*{Deux ouverts et le carré de Mayer--Vietoris (pages 42 à 45)}

Se donner deux ouverts $U, U'$ de $X$ revient à se donner
$\varphi : X \to Q$, avec $U = \varphi^{-1}(\{0\} \times \Delta^1)$ et
$U' = \varphi^{-1}(\Delta^1 \times \{0\})$. Ses fibres sont $X_{0,0} = U \cap
U'$, $X_{0,1} = U \cap T'$, $X_{1,0} = T \cap U'$ et $X_{1,1} = T \cap T'$,
où $T$ et $T'$ sont les complémentaires de $U$ et $U'$. La condition
$U \cup U' = X$ équivaut à $X_{1,1} = \varnothing$, donc à ce que $\varphi$ se
factorise par $\Psi$. Les fibres sont alors
$X_0 = U \cap U'$, $X_a = U \smallsetminus U'$, $X_b = U' \smallsetminus U$,
et
\[
U = X_{/a}, \qquad U' = X_{/b}, \qquad U \cap U' = X_{/0} .
\]
\footnote{Il écrit $\{1,1\}$ pour le point $(1,1)$. Les ouverts
$\{0, a\}$ et $\{0, b\}$ de $\Psi$, dont $U$ et $U'$ sont les images
réciproques, sont notés $U_0$, $U'_0$, soulignés.}

Le diagramme $U \leftarrow U \cap U' \rightarrow U'$ de $\mathrm{Cat}$ est
donc le foncteur $\Psi \to \mathrm{Cat}$, $s \mapsto X_{/s}$, et il définit
une catégorie cofibrée scindée $\overline{X} = \mathcal{G}_\Psi(X)$ sur
$\Psi$, de fibres $\overline{X}_a = U$, $\overline{X}_b = U'$,
$\overline{X}_0 = U \cap U'$, les foncteurs de transition étant les
inclusions. Le foncteur $\alpha : X \to \overline{X}$, $x \mapsto
(\varphi(x), x)$, est dans $W$.\footnote{La page écrit
« $\alpha \in W_{\Psi}$ », l'indice surchargé et lu avec doute, et
ne justifie pas. Voici pourquoi, pour toute classe $W$ qui contient les
foncteurs ayant un adjoint. Le foncteur $r : \overline{X} \to X$, $(s, x)
\mapsto x$, vérifie $r \alpha = \mathrm{id}_X$, et les flèches $(\varphi(x),
x) \to (s, x)$ forment une transformation naturelle $\alpha r \Rightarrow
\mathrm{id}$. Ainsi $\alpha$ est une équivalence d'homotopie au sens des
catégories, et est dans $W_\infty$, comme dans $W_{\mathbb D}$. L'argument est
le nôtre. Les deux traits vers $\Psi$ du triangle sont sans pointe sur la
page.}

La projection $\overline{\varphi} : \overline{X} \to \Psi$ est propre, de
sorte que $\overline{\varphi}_*$ se calcule fibre par fibre. Pour un
coefficient constant $\xi_{\overline{X}}$, $\xi \in \mathcal{A}$, la valeur de
$\overline{\varphi}_*(\xi_{\overline{X}})$ en $0$, $a$, $b$ est
$H^\bullet_{\mathbb D}(U \cap U', \xi)$, $H^\bullet_{\mathbb D}(U, \xi)$,
$H^\bullet_{\mathbb D}(U', \xi)$ ; le diagramme
$H^\bullet_{\mathbb D}(U, \xi) \to H^\bullet_{\mathbb D}(U \cap U', \xi)
\leftarrow H^\bullet_{\mathbb D}(U', \xi)$, de type $\Psi^{\mathrm{op}}$, est
donc le diagramme sous-jacent de
\[
\gamma_0 = \overline{\varphi}_*(\xi_{\overline{X}}) \in \mathbb D(\Psi).
\]
Soit $\gamma = k_*(\gamma_0) \in \mathbb D(Q)$. Un objet de $\mathbb D(Q)$
de la forme $k_*(\gamma_0)$ est ce qu'on appelle un carré
\emph{cartésien} ; comme $k$ est pleinement fidèle, $k^*\gamma \simeq
\gamma_0$.\footnote{La définition, qui est la reçue, n'est pas écrite sur la
page ; elle y dit « le “carré cartésien” », entre guillemets.} Le carré
sous-jacent à $\gamma$ est

\begin{tikzcd}[column sep=small]
& H^{\bullet}_{\mathbb{D}}(X, \xi) \arrow[dl, "i^*"'] \arrow[dr, "i'^*"] & \\
H^{\bullet}_{\mathbb{D}}(U, \xi) \arrow[dr, "i_0'^*"'] & & H^{\bullet}_{\mathbb{D}}(U', \xi) \arrow[dl, "i_0^*"] \\
& H^{\bullet}_{\mathbb{D}}(U \cap U', \xi) &
\end{tikzcd}

où $i : U \to X$, $i' : U' \to X$, $i'_0 : U \cap U' \to U$ et
$i_0 : U \cap U' \to U'$ sont les inclusions. En effet, si
$\psi = k \circ \overline{\varphi} : \overline{X} \to Q$, alors $\gamma =
\psi_*(\xi_{\overline{X}})$, d'où
\[
\gamma(1,1) \simeq H^\bullet_{\mathbb D}(Q, \gamma) \simeq
H^\bullet_{\mathbb D}(\overline{X}, \xi) \simeq H^\bullet_{\mathbb D}(X, \xi),
\]
la première identification parce que $(1,1)$ est final dans $Q$, la dernière
parce que $\alpha \in W$.\footnote{En marge, entre les deux moitiés de la
feuille : « à vérifier (ou à mettre comme axiome sur les dérivateurs) ». Une
note encadrée ajoute que le carré est cartésien pour tout
$\xi \in \mathbb D(X)$, et pas seulement pour $\xi$ constant.}

Dans un carré cartésien, l'image réciproque d'un isomorphisme est un
isomorphisme : si $i_0^*$ est un isomorphisme, $i^*$ en est un. Ceci valant
pour tout $\xi \in \mathcal{A}$,
\[
i_0 \in W_{\mathbb D} \Longrightarrow i \in W_{\mathbb D} .
\]
\footnote{La page écrit « $i_0^*$ iso $\Longleftrightarrow i^*$ iso » comme
une propriété générale des carrés cartésiens dans $\mathcal{A}$. Seule
l'implication $\Rightarrow$ est générale ; la réciproque demande que le carré
soit aussi cocartésien, comme la page 46 le dit elle-même. La
conclusion qu'elle tire, $i_0 \in W_{\mathbb D} \Rightarrow i \in
W_{\mathbb D}$, n'utilise que $\Rightarrow$.} C'est la moitié de W7 : pour
deux ouverts $U, U'$ recouvrant $X$, si $U \cap U' \to U'$ est une
équivalence, $U \to X$ en est une.

\pagerange{46}{48}
\subsection*{La réciproque, et le « Principe » (pages 46 à 48)}

Pour avoir la réciproque $i \in W_{\mathbb D} \Rightarrow i_0 \in
W_{\mathbb D}$, il faudrait savoir que le carré $\gamma$ est non seulement
cartésien, mais aussi \emph{cocartésien}. « Moralement », qu'il soit cartésien
lui semble la traduction cohomologique de ce que le carré de cribles
$U \cap U' \to U, U' \to X$ est cocartésien dans $\mathrm{Cat}$, en un sens
assez fort. Ce carré est aussi trivialement cartésien, et l'on peut s'attendre
à ce que son transformé par $H^\bullet_{\mathbb D}(-, \xi)$ soit
cocartésien.\footnote{Quatre lignes biffées ouvrent la page 46 ; on y lit
« Dérivateur » avec doute.}

Il se place alors sur $Q$ plutôt que sur $\Psi$. On a $\mathcal{G}_Q(X)|_\Psi
\simeq \mathcal{G}_\Psi(X)$, et la quatrième fibre est $X_{/(1,1)} \simeq X$,
de sorte que le diagramme de type $Q$ de $\mathrm{Cat}$ qui définit
$\mathcal{G}_Q(X)$ est le carré des cribles lui-même. Ce carré n'est pas
quelconque. Sa propriété essentielle est que, pour $V = U \cap U'$, le
foncteur
\[
\mathcal{G}_\Psi(X) = \int \bigl(V \to U,\ V \to U'\bigr) \longrightarrow X
\]
est dans $W$.\footnote{C'est le foncteur $r$ de la note précédente. Par le
théorème de Thomason (1979), la construction de Grothendieck d'un diagramme de
catégories en calcule la limite inductive homotopique ; dire que ce foncteur
est une équivalence, c'est dire que le carré des cribles est
homotopiquement cocartésien. La glose est la nôtre. Sous l'intégrale, la
source commune est d'abord écrite puis surchargée ; on la lit $V$.}

\textbf{« Principe » (?).} Soit $X_a \leftarrow X_0 \rightarrow X_b$ un
diagramme de $\mathrm{Cat}$ de type $\Psi$, $\mathcal{X}$ la catégorie
cofibrée sur $\Psi$ qu'il définit et $\varphi : \mathcal{X} \to \Psi$ la
projection. Pour $\xi \in \mathbb D(\mathcal{X})$, soit $\gamma_0 =
\varphi_*(\xi) \in \mathbb D(\Psi)$ et $\gamma = k_*(\gamma_0) \in \mathbb
D(Q)$. Alors
\[
\gamma(0,0) \simeq H^\bullet_{\mathbb D}(X_0, \xi), \quad
\gamma(0,1) \simeq H^\bullet_{\mathbb D}(X_a, \xi), \quad
\gamma(1,0) \simeq H^\bullet_{\mathbb D}(X_b, \xi), \quad
\gamma(1,1) \simeq H^\bullet_{\mathbb D}(\mathcal{X}, \xi),
\]
et le carré sous-jacent est celui des restrictions. Le carré
$X_0 \to X_a, X_b \to \mathcal{X}$ des fibres n'est même pas commutatif, car
$X_a \cap X_b = \varnothing$ dans $\mathcal{X}$ ; mais celui des
$X_{/0} \subset X_{/a}, X_{/b} \subset \mathcal{X}$ l'est, et le carré des
cohomologies s'en déduit aussi.\footnote{« N'est même pas » est lu avec
doute. Le passage enjambe les deux moitiés de la feuille, que la page 48 porte
à l'italienne comme les pages 45 et 51.}

Le carré $\gamma$ est cartésien par construction, et il le serait pour tout
$\mathcal{X}$ propre sur $\Psi$, sans le supposer cofibré. Que $\mathcal{X}$
soit cofibré doit donner en plus qu'il est \emph{cocartésien}. Il a
l'impression que cela ne découle pas des axiomes Der 1 à Der 5, 5', et le pose
comme axiome.\footnote{La phrase est criblée de mots illisibles et de lectures
douteuses ; on n'en garde que le sens qu'impose la suite.}

\textbf{« Der 6 » (provisoire), « axiome des cofibrations ».} Si
$\mathcal{X}$ est cofibré sur $\Psi$ et $\xi \in \mathbb D(\mathcal{X})$,
l'objet $\psi_*(\xi) \in \mathbb D(Q)$, où $\psi = k \circ \varphi$, est non
seulement cartésien, mais aussi cocartésien.\footnote{« Quelconque » après
$\xi$ est biffé. Le nom « Axiome des cofibrations » est écrit en oblique au
bas de la moitié droite, souligné ; trois lignes plus petites au-dessous, et
une note oblique de la marge gauche reliée à « Der 6 », restent illisibles.}

\pagerange{49}{50}
\subsection*{Der 6', et la forme symétrique (pages 49 et 50)}

\textbf{« Der 6' » (provisoire).} Si $\mathcal{X}$ est fibrée sur
$\Phi = \Psi^{\mathrm{op}} = (a \to 0 \leftarrow b)$ (donc donnée, à
$\Phi$-équivalence près, par un diagramme $X_a \leftarrow X_0 \rightarrow X_b$
de $\mathrm{Cat}$) et $\xi \in \mathbb D(\mathcal{X})$ est constant, alors
$\psi_!(\xi)$, pour $\psi : \mathcal{X} \to \Phi \subset Q^{\mathrm{op}}$, est
non seulement cocartésien, « sans mérite », mais aussi
cartésien.\footnote{La page appelle $1$ le sommet de $\Phi$, « $(a \to 1
\leftarrow b)$ » ; il deviendra $0$ aux pages 51 et 52, et on l'appelle $0$
partout. Elle écrit $\psi_!(\xi) \in \mathbb D(Q)$, mais dessine
$\Phi$ dans $Q^\circ$, ses trois sommets entourés ; c'est
$\mathbb D(Q^{\mathrm{op}})$. Le diagramme définissant $\mathcal{X}$ est
reproduit tel que la page le trace (un trait sans pointe entre $X_a$ et $X_0$,
une flèche de $X_a$ vers $X_b$) ; la forme d'un diagramme de fibres d'une
catégorie fibrée sur une co-étoile est celle qu'on écrit. Dans la marge
gauche, en oblique, « Axiome des fibrations », lu avec doute.}

\emph{NB.} Der 6' pour $\mathbb D$ n'est autre que Der 6 pour le dérivateur
opposé $\mathbb D^\circ$, $\mathbb D^\circ(I) = \mathbb
D(I^{\mathrm{op}})^{\mathrm{op}}$. Il ne voit pas comment l'un impliquerait
l'autre, hors le cas exceptionnel où il y a un foncteur dualisant dans
$\mathcal{A}$, et une dualité correspondante entre cohomologie et homologie.

Il semble finalement que l'axiome qui englobe les deux soit le suivant :

\textbf{« Der 6 ».} Un objet $\gamma \in \mathbb D(\Delta^1 \times
\Delta^1)$ est cartésien si et seulement s'il est cocartésien.

C'est, pour un dérivateur pointé, la définition reçue aujourd'hui
d'un dérivateur \emph{stable}.\footnote{Sous cette forme : Maltsiniotis et
Groth, cités de mémoire. Les catégories $\mathbb D(I)$ d'un dérivateur stable
sont triangulées. Le dérivateur des types d'homotopie ne l'est pas : un carré
homotopiquement cocartésien d'espaces n'est pas en général cartésien (voir la
note de la page 51). « Der 6 » est suivi d'un mot noirci ; l'addition entre
crochets, interlinéaire et très raturée, et la note oblique de la marge
gauche (« l'axiome “vrai” … cartésien »), ne se lisent que par fragments.}

Page 50, sur le carré

\begin{tikzcd}[column sep=small, row sep=small]
& \gamma(0,1) \arrow[dl, "j_0"'] & \\
\gamma(0,0) & & \gamma(1,1) \arrow[ul, "i_1"'] \arrow[dl, "j_1"] \\
& \gamma(1,0) \arrow[ul, "i_0"] &
\end{tikzcd}

un carré bicartésien a $i_0$ iso si et seulement si $i_1$ l'est, et $j_0$ iso
si et seulement si $j_1$ l'est : « cartésien » donne un sens, « cocartésien »
l'autre.\footnote{La page n'écrit que « $i_1$ iso $\Rightarrow i_0$ iso, $j_1$
iso $\Rightarrow j_0$ iso », qui est le sens cocartésien ($i_0$ est le
changement de cobase de $i_1$), puis la parenthèse « car cart.\ implique
$\Longrightarrow$, cocartésien implique $\Longleftarrow$ ». En marge, à côté du
carré : il faudra invoquer aussi les dérivateurs induits $\mathbb D_I$
($I \in \mathrm{Diag}$), dans $\mathbb D(I \times \Delta^1 \times \Delta^1)$.}
La dernière forme se déduirait aussitôt de la précédente, pourvu qu'on ait le
lemme suivant.\footnote{« Montre aussitôt » est lu avec doute, et la suite de
la phrase illisible ; le sens de la déduction n'est pas écrit au-delà.}

\textbf{« Lemme ».} Soit $\xi = (\xi_a \xrightarrow{j_0} \xi_0
\xleftarrow{i_0} \xi_b) \in \mathbb D(\Psi)$ tel que $i_0$ soit un
isomorphisme. Alors $\xi$ est isomorphe à l'image réciproque $\rho^*\eta$
d'un $\eta \in \mathbb D(\Psi_a)$, où $\Psi_a = (0 \to a)
\simeq \Delta^1$ est l'ouvert $\{0, a\}$ de $\Psi$ et $\rho : \Psi \to
\Psi_a$ la rétraction qui envoie $b$ sur $0$.

La page veut montrer que la flèche d'adjonction $\rho^*\rho_*(\xi)
\to \xi$ est un isomorphisme. Comme $\rho$ est lisse, induite par la première
projection $\Delta^1 \times \Delta^1 \to \Delta^1$, $\rho_*$ se calcule fibre
par fibre, et l'on est ramené à ce que $\mathcal{A} \to \mathbb
D^{\mathrm{cc}}(\Delta^1)$ soit une équivalence, « forme de Der 3 ». Elle
s'arrête là.\footnote{L'exposant « cc » se lit mal ; une note verticale biffée
le long de la marge gauche reste illisible. Le lemme est vrai, et se démontre
plus directement. Soit $j : \Psi_a \subset \Psi$ ; on a $\rho j =
\mathrm{id}$ et une transformation naturelle $j\rho \Rightarrow
\mathrm{id}_\Psi$ (car $j\rho(b) = 0 \leq b$). Elle induit une flèche entre
$\rho^* j^*\xi$ et $\xi$ dont les composantes en $0$ et $a$ sont des identités
et celle en $b$ est $i_0$. C'est donc un isomorphisme, les isomorphismes se
lisant sur les valeurs, et l'on prend $\eta = j^*\xi$. Cette démonstration est
la nôtre.}

\pagerange{51}{52}
\subsection*{HOT, Der 8, W7 bis et W8 bis (pages 51 et 52)}

« Dans HOT, j'ai admis que : tout carré cocartésien est cartésien. Pour qu'un
carré cartésien

\begin{tikzcd}[column sep=small, row sep=small]
& \xi_0 \arrow[dl] \arrow[dr] & \\
\xi_a \arrow[dr] & & \xi_b \arrow[dl] \\
& \xi_1 &
\end{tikzcd}

soit aussi cocartésien, il faut et il suffit que l'on ait
$\pi_0(\xi_1) = \mathrm{Im}\, \pi_0(\xi_a) \cup \mathrm{Im}\, \pi_0(\xi_b)$. »
\footnote{HOT n'est pas identifié par le dossier. Les deux énoncés valent pour
les complexes connectifs (concentrés en degrés homologiques positifs), où
$\pi_0$ est $H_0$ et la réunion des images une somme. Une somme amalgamée
homotopique d'objets connectifs y est encore connective, donc cartésienne. Un
produit fibré homotopique y est connectif, donc cocartésien, exactement quand
$H_0(\xi_a) \oplus H_0(\xi_b) \to H_0(\xi_1)$ est surjective. Ils sont faux
pour les types d'homotopie d'espaces : la somme amalgamée de $\mathrm{pt}
\leftarrow S^0 \to \mathrm{pt}$ est $S^1$, et le produit fibré homotopique de
$\mathrm{pt} \to S^1 \leftarrow \mathrm{pt}$ est $\Omega S^1 \simeq \mathbb Z$,
non $S^0$. La remarque est la nôtre. Les trois lignes sont tenues par un trait
vertical à gauche ; une note oblique de la marge, en partie lisible, porte
« $\xi_1 = e$ ». La ligne suivante, « Der 6 ou Der 6' … donc … », est
illisible pour l'essentiel.}

W8 semble découler d'un principe général :

\textbf{« Der 8 ».} Soit $f : \xi \to \eta$ un morphisme de carrés de
$\mathcal{A}(\Delta^1 \times \Delta^1)$, de composantes $f_0, f_a, f_b,
f_1$, et les conditions a) $f_0$ est iso ; b) $f_1$ est iso ; c) $f_a$ et
$f_b$ sont iso. Si les carrés sont cartésiens, b) et c) entraînent a) ; s'ils
sont cocartésiens, a) et c) entraînent b) ; s'ils sont bicartésiens, de plus a)
et b) entraînent c) (?).\footnote{Le cube est redessiné ; sur la page les deux carrés sont décalés
l'un sous l'autre, les quatre flèches verticales se croisent et les indices
$a$, $b$ se lisent mal. « Cartésiens », « cocartésiens », « bicartésiens »
sont lus avec doute ; un \emph{NB} qui suit est barré. Les deux premières
implications sont celles de l'image réciproque et de l'image directe d'un
isomorphisme. La troisième, marquée « (?) », vaut dans un dérivateur stable,
où un carré bicartésien donne une suite exacte longue de Mayer--Vietoris
$\cdots \to \xi_0 \to \xi_a \oplus \xi_b \to \xi_1 \to \cdots$, et le lemme des
cinq rend $f_a \oplus f_b$ inversible, donc $f_a$ et $f_b$. La remarque est la
nôtre ; en marge gauche : « Il suffit d'exiger $a + b \Rightarrow c$ … ».}

\begin{tikzcd}[column sep=small, row sep=small]
& \xi_a \arrow[dr] \arrow[ddd, "f_a"] & \\
\xi_0 \arrow[ur] \arrow[dr] \arrow[ddd, "f_0"'] & & \xi_1 \arrow[ddd, "f_1"] \\
& \xi_b \arrow[ur] \arrow[ddd, "f_b"'] & \\
& \eta_a \arrow[dr] & \\
\eta_0 \arrow[ur] \arrow[dr] & & \eta_1 \\
& \eta_b \arrow[ur] &
\end{tikzcd}


\textbf{W7 bis.} Soit $X_a \leftarrow X_0 \rightarrow X_b$ un diagramme de
$\mathrm{Cat}$, et $\mathcal{X}$ la catégorie fibrée sur $\Phi = (a \to 0
\leftarrow b)$ qu'il définit. On a un diagramme commutatif dans
$\mathrm{Hot}_W = W^{-1}\mathrm{Cat}$ :

\begin{tikzcd}[column sep=small, row sep=small]
& X_0 \arrow[dl, "i_a"'] \arrow[dr, "i_b"] & \\
X_a \arrow[dr, "i'_b"'] & & X_b \arrow[dl, "i'_a"] \\
& \mathcal{X} &
\end{tikzcd}

(dans $\mathrm{Cat}$, les deux composés sont homotopes par un chemin de
longueur 2, homotopes tous deux, par une homotopie de longueur 1, à
l'inclusion $X_0 \hookrightarrow \mathcal{X}$). Alors $i_a \in W$ si et
seulement si $i'_a \in W$, et $i_b \in W$ si et seulement si $i'_b \in
W$.\footnote{Les pointes des flèches de $\Phi$ se lisent mal, ici et page 52,
où la page écrit une fois « $(0 \leftarrow a,\ 0 \to b)$ » ; on écrit
$\Phi = (a \to 0 \leftarrow b)$, la forme sur laquelle une catégorie fibrée
a pour fibres un diagramme $X_a \leftarrow X_0 \rightarrow X_b$, et celle de
la page 49. En marge : « Autre possibilité : $\mathcal{X}$ fibrée … ».}

\textbf{W8 bis.} Soit un cube commutatif de $\mathrm{Cat}$ formé de deux tels
diagrammes et de $f_0, f_a, f_b$, complété dans $\mathrm{Hot}_W$ par
$f : \mathcal{X} \to \mathcal{X}'$, où $\mathcal{X} = \int(X_0 \to X_a,\ X_0
\to X_b)$ et de même $\mathcal{X}'$. Considérons les conditions a) $f \in W$ ;
b) $f_0 \in W$ ; c) $f_a$ et $f_b \in W$. Deux d'entre elles entraînent la
troisième.\footnote{Cube redessiné sur le modèle du précédent ; les flèches vers
$\mathcal{X}$ et $\mathcal{X}'$ sont pointillées sur la page, les autres
pleines, et les étiquettes $f_a$, $f_b$ se lisent mal. Deux notes obliques de
part et d'autre du cube, en partie biffées, ne se lisent que par fragments.}

\begin{tikzcd}[column sep=small, row sep=small]
& X_a \arrow[dr, dashed] \arrow[ddd, "f_a"] & \\
X_0 \arrow[ur] \arrow[dr] \arrow[ddd, "f_0"'] & & \mathcal{X} \arrow[ddd, "f"] \\
& X_b \arrow[ur, dashed] \arrow[ddd, "f_b"'] & \\
& X'_a \arrow[dr, dashed] & \\
X'_0 \arrow[ur] \arrow[dr] & & \mathcal{X}' \\
& X'_b \arrow[ur, dashed] &
\end{tikzcd}


La page 52 (sa page 18, numéro entouré) en donne une mise au propre, sous le
titre « Corollaire de W7 bis », biffé en « Variante forte ». \emph{W7 bis} :
soit $\mathcal{X}$ lisse sur $\Phi$, d'où le même diagramme ; alors
$i_a \in W$ si et seulement si $i'_a \in W$ (donc de même pour $i_b$,
$i'_b$). Cela implique l'axiome dual W7' bis, pour $\mathcal{X}$ propre sur
$\Psi$. \emph{W8 bis} : soient $\mathcal{X}$, $\mathcal{X}'$ lisses sur
$\Phi$ et $f : \mathcal{X} \to \mathcal{X}'$ un $\Phi$-morphisme, induisant
$f_0, f_a, f_b$ sur les fibres ; deux des conditions a) $f \in W$, b) $f_0 \in
W$, c) $f_a, f_b \in W$ entraînent la troisième.\footnote{« Lisse »
généralise « fibré » comme « propre » généralise « cofibré », selon la
convention de la page 44.}

Ces deux énoncés ont un contenu qu'on peut dire, bien que la page ne le dise
pas. La catégorie $\mathcal{X}$ est la somme amalgamée homotopique du
diagramme. Les implications « $i_a \in W \Rightarrow i'_a \in W$ » et « b) et
c) $\Rightarrow$ a) » expriment l'invariance homotopique de cette somme : elles
valent pour $W_\infty$ comme pour toute classe $W_{\mathbb D}$. Les
implications inverses sont d'une autre nature. Elles sont fausses pour
$W_\infty$ et vraies pour les équivalences d'homologie, par la suite de
Mayer--Vietoris et le lemme des cinq, comme pour $W_{\mathbb D}$ quand
$\mathbb D$ est stable.\footnote{Remarque et contre-exemple sont les nôtres.
Soit $X_0$ une catégorie dont le nerf est un espace acyclique non contractile,
comme le complémentaire d'un point dans la sphère d'homologie de Poincaré ;
soit $X_a = e$ et $X_b$ une catégorie contractile contenant $X_0$, par exemple
$X_0$ muni d'un objet final ajouté. Alors $\mathcal{X}$ a le type de la
suspension de $X_0$, simplement connexe et acyclique, donc contractile. Ainsi
$i'_b : X_a \to \mathcal{X}$ est dans $W_\infty$ sans que $i_b : X_0 \to X_b$
le soit. De même, pour W8 bis, l'application de ce diagramme vers $e
\leftarrow e \rightarrow e$ satisfait a) et c), pas b).} C'est ce qui fait de
« Der 6 » l'axiome que ces pages cherchent : pour qu'un recollement se lise
dans les deux sens, il faut être dans le monde stable, où cartésien et
cocartésien coïncident.

\end{document}
